\documentclass[12pt,a4paper,UTF8]{ctexart}
\usepackage[a4paper,top=22mm,bottom=22mm,left=22mm,right=22mm,headheight=15pt]{geometry}
\usepackage{amsmath,amssymb,mathtools,bm}
\usepackage{graphicx,xcolor,booktabs,longtable,array,enumitem,fancyhdr,caption,float,microtype,setspace}
\newcommand{\Needspace}[1]{\par\begingroup\dimen0=#1\relax\dimen2=\pagegoal\advance\dimen2 by-\pagetotal\ifdim\dimen0>\dimen2\newpage\fi\endgroup}
\newcommand{\needspace}[1]{\par\begingroup\skip0=\lastskip\vskip-\skip0\dimen0=#1\relax\advance\dimen0 by\skip0\vskip\dimen0\penalty-100\vskip-\dimen0\vskip\skip0\penalty9999\endgroup}
\usepackage[hidelinks,unicode]{hyperref}
\usepackage[most]{tcolorbox}
\definecolor{ContestBlue}{HTML}{17365D}
\hypersetup{pdftitle={第 050 套},pdfauthor={}}
\setstretch{1.24}
\setlength{\parindent}{2em}
\setlength{\parskip}{0.35em}
\setlength{\emergencystretch}{2em}
\setlist{itemsep=0.35em,topsep=0.4em}
\captionsetup{font=small,labelfont=bf,justification=centering}
\setcounter{secnumdepth}{0}
\pagestyle{fancy}
\fancyhf{}
\fancyhead[L]{\small 第 050 套}
\fancyfoot[C]{\thepage}
\renewcommand{\headrulewidth}{0.4pt}
\ctexset{
  section={format=\Large\bfseries\color{ContestBlue},beforeskip=1.2em,afterskip=0.7em},
  subsection={format=\large\bfseries\color{ContestBlue},beforeskip=1.0em,afterskip=0.5em},
  subsubsection={format=\normalsize\bfseries,beforeskip=0.8em,afterskip=0.35em}
}
\makeatletter
\newcommand{\tightlist}{\setlength{\itemsep}{0pt}\setlength{\parskip}{0pt}}
\newcommand{\maxwidth}{\ifdim\Gin@nat@width>\linewidth\linewidth\else\Gin@nat@width\fi}
\newcommand{\maxheight}{\ifdim\Gin@nat@height>0.72\textheight 0.72\textheight\else\Gin@nat@height\fi}
\newcommand{\pandocbounded}[1]{\begingroup\setkeys{Gin}{width=\maxwidth,height=\maxheight,keepaspectratio}#1\endgroup}
\makeatother
\begin{document}
\begin{titlepage}
  \centering
  \vspace*{58mm}
  {\Huge\bfseries\color{ContestBlue} 第 050 套\par}
  \vfill
  \vspace*{28mm}
\end{titlepage}
\section{第 1 题}\label{ux7b2c-1-ux9898}

空气以速度 \(v_0\) 竖直向下流向挡板，其中比例为 \(\eta=0.1\)
的空气分子被挡板近似原速弹回。把颗粒物先近似为直径和高度均为
\(2r\)、密度为 \(\rho\)
的均匀圆柱，圆柱轴线始终竖直；空气和颗粒在进入挡板附近时均以 \(v_0\)
向下运动。反弹分子与颗粒发生完全弹性碰撞，且只有在颗粒距挡板不超过
\(h_0\)
时碰撞作用才显著。忽略颗粒与其他空气分子、颗粒之间的相互作用。已知 \[
v_0\sim10\,\mathrm{m\,s^{-1}},\qquad h_0=1.0\ {\rm cm},\qquad
\rho_0=1.29\ {\rm kg\,m^{-3}},\qquad
\rho=2.0\times10^3\ {\rm kg\,m^{-3}}.
\]

（1）估算空气分子碰撞使圆柱颗粒获得的加速度数量级；

（2）求能够随气流离开的圆柱颗粒的最大半径 \(r_0\)；

（3）若颗粒改为半径为 \(r\)
的球，空气分子按镜面反射规律与球面碰撞，求最大半径 \(r_0'\)。

\begin{figure}[H]
\centering
\includegraphics[width=0.778\linewidth,height=\textheight,keepaspectratio,alt={图1-1　颗粒与挡板反射气流的相对运动}]{../figures/question-01.pdf}
\caption*{图1-1　颗粒与挡板反射气流的相对运动}
\end{figure}

\subsection{一、题型定位}\label{ux4e00ux9898ux578bux5b9aux4f4d}

颗粒所受冲击阻力由空气与颗粒的相对速度决定。颗粒减速时这个速度也会变化；恒定初始阻力只适合做估算，不能代替完整的停止距离。

\subsection{二、知识点}\label{ux4e8cux77e5ux8bc6ux70b9}

\begin{enumerate}
\def\labelenumi{\arabic{enumi}.}
\tightlist
\item
  相对速度决定碰撞通量；
\item
  完全弹性正碰的动量改变量；
\item
  质量流率乘以单位质量冲量得到平均力；
\item
  匀减速停止距离 \(s=v^2/(2a)\)；
\item
  球面元 \(dA=2\pi r^2\sin\theta\,d\theta\) 的投影与法向分解；
\item
  数量级、量纲和极限检验。
\end{enumerate}

\subsection{三、物理图像}\label{ux4e09ux7269ux7406ux56feux50cf}

挡板附近向上的反射分子流与仍向下运动的颗粒迎面相遇。两者相对速率不是
\(v_0\)，而是 \[
U=v_0-(-v_0)=2v_0.
\]
圆柱底面处处具有同一法向，所有碰撞的竖直冲量可以直接相加。球面上越靠近侧面的面元，迎风投影越小，法向冲量在竖直方向的分量也越小，所以球受到的总力比同半径圆柱小。

\subsection{四、解题入口}\label{ux56dbux89e3ux9898ux5165ux53e3}

对面积为 \(A\) 的迎风面，在时间 \(dt\) 内扫过的相对体积为
\(AUdt\)。真正反射并参与碰撞的空气质量为 \[
dm_{\rm air}=\eta\rho_0AUdt.
\]
完全弹性反射使这部分空气相对于颗粒的法向速度反向，颗粒得到的冲量大小为
\(2U\,dm_{\rm air}\)。因此平面迎风面的平均力为 \[
F=2\eta\rho_0AU^2.
\] 这条``质量通量 \(\times\) 单位质量动量变化''的式子是全题入口。

\subsection{五、解析}\label{ux4e94ux89e3ux6790}

\needspace{3\baselineskip}
\textbf{（一）圆柱颗粒所受平均力}

圆柱底面积为 \[
A=\pi r^2.
\] 相对速度 \(U=2v_0\)，故单位时间参与碰撞的空气质量为 \[
\frac{dm_{\rm air}}{dt}
=\eta\rho_0\pi r^2U.
\] 每千克空气在完全弹性正碰中法向速度由 \(+U\) 变为
\(-U\)，动量改变量的大小为
\(2U\)。由牛顿第三定律，颗粒受到大小相等、方向向上的冲量，所以 \[
\begin{aligned}
F_{\rm cyl}
&=2U\frac{dm_{\rm air}}{dt}\\
&=2\eta\rho_0\pi r^2U^2\\
&=2\eta\rho_0\pi r^2(2v_0)^2\\
&=8\eta\rho_0\pi r^2v_0^2.
\end{aligned} \tag{1}
\]

圆柱体积和质量分别是 \[
V_{\rm cyl}=\pi r^2(2r)=2\pi r^3,
\qquad
m_{\rm cyl}=2\pi\rho r^3.
\] 所以向上加速度为 \[
\begin{aligned}
a_{\rm cyl}
&=\frac{F_{\rm cyl}}{m_{\rm cyl}}\\
&=\frac{8\eta\rho_0\pi r^2v_0^2}{2\pi\rho r^3}\\
&=\boxed{\frac{4\eta\rho_0v_0^2}{\rho r}}.
\end{aligned} \tag{2}
\] 式（2）说明半径越小，加速度越大，这是过滤装置能按粒径筛选的根本原因。

\begin{figure}[H]
\centering
\includegraphics[width=0.82\linewidth,height=\textheight,keepaspectratio,alt={图1-2　圆柱与球所受反射气流冲量的比较}]{../figures/solution-01.pdf}
\caption*{图1-2　圆柱与球所受反射气流冲量的比较}
\end{figure}

\needspace{3\baselineskip}
\textbf{（二）最大圆柱半径}

颗粒进入有效碰撞层时仍以 \(v_0\)
向下运动。要在撞到挡板以前停下，所需停止距离不得大于 \(h_0\)： \[
s_{\rm stop}=\frac{v_0^2}{2a_{\rm cyl}}\le h_0. \tag{3}
\] 临界颗粒恰好在距离 \(h_0\) 内停下。将式（2）代入等号条件： \[
\begin{aligned}
\frac{v_0^2}{2[4\eta\rho_0v_0^2/(\rho r_0)]}
&=h_0,\\
\frac{\rho r_0}{8\eta\rho_0}&=h_0,\\
r_0&=\frac{8\eta\rho_0h_0}{\rho}.
\end{aligned}
\] 代入数值： \[
\begin{aligned}
r_0
&=\frac{8(0.1)(1.29)(1.0\times10^{-2})}{2.0\times10^3}\ {\rm m}\\
&=5.16\times10^{-6}\ {\rm m}.
\end{aligned}
\] 因此 \[
\boxed{r_0\simeq5.2\ \mu{\rm m}}. \tag{4}
\] 在临界状态下 \[
a_{\rm crit}=\frac{v_0^2}{2h_0}.
\] 取题中所说 \(v_0\sim10\ {\rm m\,s^{-1}}\)，得到 \[
a_{\rm crit}\sim\frac{10^2}{2\times10^{-2}}
=5\times10^3\ {\rm m\,s^{-2}}
\sim10^4\ {\rm m\,s^{-2}}. \tag{5}
\] 与重力加速度相比大约高三个数量级，因此前面的主计算忽略重力是自洽的。

\needspace{3\baselineskip}
\textbf{（三）球形颗粒的球面积分}

在球面迎风半球上，以最下端为极点，用 \(\theta\)
表示面元法线与竖直方向的夹角。环带面积为 \[
dA=2\pi r^2\sin\theta\,d\theta.
\] 相对气流穿过该面元的法向速度是 \(U\cos\theta\)，所以撞击质量流率为 \[
d\dot m=\eta\rho_0U\cos\theta\,dA.
\] 每单位质量的法向冲量是 \(2U\cos\theta\)，该冲量在竖直方向还要再乘
\(\cos\theta\)。于是 \[
\begin{aligned}
dF_z
&=(d\dot m)(2U\cos\theta)\cos\theta\\
&=2\eta\rho_0U^2\cos^3\theta\,dA\\
&=4\pi\eta\rho_0r^2U^2
\sin\theta\cos^3\theta\,d\theta.
\end{aligned}
\] 对迎风半球积分： \[
\begin{aligned}
F_{\rm sph}
&=4\pi\eta\rho_0r^2U^2
\int_0^{\pi/2}\sin\theta\cos^3\theta\,d\theta\\
&=4\pi\eta\rho_0r^2U^2
\left[-\frac{\cos^4\theta}{4}\right]_0^{\pi/2}\\
&=\eta\rho_0\pi r^2U^2\\
&=4\eta\rho_0\pi r^2v_0^2.
\end{aligned} \tag{6}
\] 球的质量为 \[
m_{\rm sph}=\frac43\pi\rho r^3,
\] 从而 \[
\begin{aligned}
a_{\rm sph}
&=\frac{F_{\rm sph}}{m_{\rm sph}}\\
&=\frac{4\eta\rho_0\pi r^2v_0^2}{(4/3)\pi\rho r^3}\\
&=\frac{3\eta\rho_0v_0^2}{\rho r}.
\end{aligned} \tag{7}
\] 再用停止距离的临界条件： \[
\frac{v_0^2}{2a_{\rm sph}}=h_0,
\] 逐步整理得 \[
\boxed{r_0'=\frac{6\eta\rho_0h_0}{\rho}
=3.87\times10^{-6}\ {\rm m}
=\frac34r_0}. \tag{8}
\]

\textbf{把估算系数与变速模型分开。} 完全弹性正碰的冲量系数确定为
\(2\)，不能把漏掉该因子说成另一种约定。正文把相对速率取为入层时的
\(2v_0\)，所以 \(5.16\,\mu\mathrm m\)
是常加速度估算。若反射流密度不变而保留颗粒减速，向下坐标 \(z\)、速率
\(v\) 满足 \[
v\frac{\mathrm dv}{\mathrm dz}
=-\frac{\eta\rho_0}{\rho r}(v_0+v)^2.
\] 从 \(v=v_0\) 积到 \(0\)，停止距离为 \[
s_{\rm stop}=\frac{\rho r}{\eta\rho_0}
\int_0^1\frac{u\,\mathrm du}{(1+u)^2}
=\frac{\rho r}{\eta\rho_0}\left(\ln2-\frac12\right).
\] 由此圆柱截止半径为
\(r_{0,\rm var}=\eta\rho_0h_0/[\rho(\ln2-1/2)]\approx3.34\,\mu\mathrm m\)，球形为其
\(3/4\approx2.50\,\mu\mathrm m\)。两者都保持正文的粒径尺度和形状比，但数值系数不能混用。估算忽略吸附质量增长、分子热运动和重力。

\textbf{按实际相对速度积分停止距离。} 取颗粒向下速度为
\(v\ge0\)，空气向上速度为 \(v_0\)，相对速率为
\(U=v_0+v\)。碰撞模型的阻力可写成 \[
F=\eta\frac{m\rho_0}{\rho r}(v_0+v)^2.
\] 其中 \(\eta\) 包含粒子形状和碰撞传递系数。由 \(m\dot v=-F\)，并用
\(\dot v=v\,dv/dz\)，得到 \[
v\frac{dv}{dz}=-\eta\frac{\rho_0}{\rho r}(v_0+v)^2.
\] 从初速 \(v_0\) 减到零时，停止距离为 \[
s=\frac{\rho r}{\eta\rho_0}\int_0^{v_0}\frac{v\,dv}{(v_0+v)^2}.
\] 令 \(u=v/v_0\)，积分不再含初速度： \[
\int_0^1\frac{u\,du}{(1+u)^2}=\ln2-\frac12.
\] 因此给定停止层厚 \(h_0\) 后，最大半径为 \[
r_{\max}=\frac{\eta\rho_0h_0}{\rho(\ln2-1/2)}.
\] 若把相对速率始终取为初值 \(2v_0\)，估计距离为
\(s_{\rm const}=\rho r/(8\eta\rho_0)\)。二者的比例是 \[
\frac{s}{s_{\rm const}}=8\ln2-4\simeq1.545.
\]
初始阻力较大，保持它不变会低估距离，从而高估可截留半径。这正解释恒定阻力估算与变速结果的差别。

球形颗粒与圆柱形颗粒的阻力面积、质量随半径的标度相同，但面积质量比不同，故在相同碰撞约定下临界半径差一个几何系数。黏附碰撞若引起颗粒质量持续增加，还须把
\(m(t)\)
写入动量方程；不能只把弹性传递系数减半后声称得到包含质量增长的完整解。

\textbf{吸附模型的两种层次。}
若暂忽略黏附空气带来的质量增长，每单位空气质量的动量变化从 \(2U\) 降为
\(U\)，圆柱减速方程变为 \[
v\frac{dv}{dz}=-\frac{\eta\rho_0}{2\rho r}(v_0+v)^2.
\] 变速停止距离为弹性反射模型的两倍： \[
s_{\rm abs}=\frac{2\rho r}{\eta\rho_0}(\ln2-1/2).
\] 同一厚度下，临界半径为 \[
r_{\rm abs,var}=\frac{\eta\rho_0h_0}{2\rho(\ln2-1/2)}\simeq1.670\,\mu\mathrm m.
\] 恒定初值力估算则给
\(2.58\,\mu\mathrm m\)。比较时必须把两套近似各自成对使用。

若真正保留质量增长，设吸附空气速度为 \(-v_0\)、颗粒向下速度为
\(v\)，则动量方程不是 \(m\dot v=F\) 的恒质量形式，而是 \[
\frac{d(mv)}{dt}=-v_0\dot m.
\] 因此 \[
m(v+v_0)=m_0(2v_0).
\] 颗粒停止时
\(m=2m_0\)，所吸附质量已经与初始质量同阶，故``质量增长可忽略''的吸附模型不自动自洽。

进一步若迎风面积 \(A\) 仍取常数，\(\dot m=\eta\rho_0A(v+v_0)\)。消去
\(v\)， \[
\dot m=\frac{2\eta\rho_0Am_0v_0}{m}.
\] 由 \(dz/dm=v/\dot m\) 积分至 \(m=2m_0\)，得到 \[
s=\frac{m_0}{4\eta\rho_0A}.
\] 对初始圆柱 \(m_0/A=2\rho r\)，临界半径成为
\(r=2\eta\rho_0h_0/\rho\)，约
\(1.29\,\mu\mathrm m\)。该结果另依赖``形状、面积不随附着改变''的假设，用来显示质量增长会改变系数，不能混称为原题的唯一吸附答案。

\textbf{停止时间与重力近似的检查。} 将变速阻力写成
\(\dot v=-a_*(v_0+v)^2\)，其中 \(a_*=\eta\rho_0/(\rho r)\)，初态
\(v=v_0\)。时间积分比路程积分少一个速度因子： \[
t_{\rm stop}=\frac1{a_*}\int_0^{v_0}\frac{dv}{(v_0+v)^2}
=\frac1{2a_*v_0}.
\] 过程中的速度可直接解出 \[
\frac1{v_0+v}=\frac1{2v_0}+a_*t.
\] 代回 \(v(t)\) 并积分位置，恰好恢复 \[
s=\frac{\ln2-1/2}{a_*}.
\] 所以停止时间随入射速度反比变化，停止距离却不含
\(v_0\)。这两个标度同时成立，是阻力随相对速度平方变化的结果。

忽略重力需要检验阻力在全程都大于重力的程度。若把重力加回，取向下为正，则
\[
\dot v=g-a_*(v_0+v)^2.
\] 减速至 \(v=0\) 的必要条件为
\(a_*v_0^2\ge g\)。若严格大于，停止发生在有限时间；若恰好等于，在
\(v\to0\) 时有 \[
\dot v\simeq-2a_*v_0v,
\] 速度指数趋零，有限时间内不会严格到零。若小于，存在向下的终端速度 \[
v_\infty=\sqrt{g/a_*}-v_0>0.
\] 因此``受到向上阻力''并不自动保证颗粒停在该气层内。

用无量纲比 \(\chi=g/(a_*v_0^2)\) 评价重力。对
\(\chi\ll1\)，全过程阻力至少约为 \(a_*v_0^2\)，忽略重力自洽；当 \(\chi\)
接近一时，停止时间和距离对参数敏感，不能继续引用无重力的数值半径作为精确阈值。

图1-2区分了空气速度、颗粒速度与相对速度。上述修正只加入重力，不改变碰撞传递系数；若又改变碰撞为吸附，还需同时使用变质量动量方程，不能把两个近似分别算出的系数随意相乘。

\newpage

\section{第 2 题}\label{ux7b2c-2-ux9898}

质量为 \(M\)、长度为 \(L\)
的均匀细杆只能在竖直平面内运动，杆的两端始终与半径为 \(R\)
的光滑半圆形碗内壁接触，且 \(R>1.5L\)。

（1）求细杆在最低平衡位置附近的微小振动周期 \(T\)；

（2）若振动的最大偏角为
\(\theta_m\)，碗对细杆的合支持力大小在转折点与通过平衡点时之差写成 \[
\Delta N=N\big|_{\theta=\theta_m}-N\big|_{\theta=0}
=\alpha Mg\theta_m^2,
\] 求 \(\alpha\)。只保留 \(\theta_m\) 的二阶小量。

\begin{figure}[H]
\centering
\includegraphics[width=0.8\linewidth,height=\textheight,keepaspectratio,alt={图2-1　细杆作为圆内定长弦的几何约束}]{../figures/question-02.pdf}
\caption*{图2-1　细杆作为圆内定长弦的几何约束}
\end{figure}

\subsection{一、题型定位}\label{ux4e00ux9898ux578bux5b9aux4f4d-1}

杆的两端同时接触圆弧，质心沿同心圆运动。把平动与绕质心转动的动能合并后，可由重力势能求周期。

\subsection{二、知识点}\label{ux4e8cux77e5ux8bc6ux70b9-1}

\begin{enumerate}
\def\labelenumi{\arabic{enumi}.}
\tightlist
\item
  圆中定长弦到圆心距离；
\item
  平行轴定理；
\item
  拉格朗日方程或定轴转动方程；
\item
  小角近似与机械能守恒；
\item
  法向、切向加速度分解；
\item
  合支持力大小的二阶展开。
\end{enumerate}

\subsection{三、物理图像}\label{ux4e09ux7269ux7406ux56feux50cf-1}

杆的两端是同一圆上的两点，因此杆始终是一条长度固定的弦。圆心到弦中点的垂线既平分弦，也垂直于杆。设杆中点即质心为
\(G\)，则 \(OG=d\) 为常量，整个杆的构型只需一个角度 \(\theta\)
描述。质心走半径为 \(d\) 的圆，杆的转角也等于 \(\theta\)。

\subsection{四、解题入口}\label{ux56dbux89e3ux9898ux5165ux53e3-1}

由直角三角形 \[
d^2+\left(\frac L2\right)^2=R^2,
\] 所以 \[
d=\sqrt{R^2-\frac{L^2}{4}}. \tag{1}
\] 然后将平动与绕质心转动合并为绕 \(O\) 的等效转动： \[
I_O=I_G+Md^2.
\]

\subsection{五、解析}\label{ux4e94ux89e3ux6790-1}

\needspace{3\baselineskip}
\textbf{（一）等效转动惯量}

均匀细杆关于质心且垂直于运动平面的轴的转动惯量是 \[
I_G=\frac{1}{12}ML^2.
\] 由平行轴定理： \[
\begin{aligned}
I_O
&=I_G+Md^2\\
&=\frac{1}{12}ML^2+M\left(R^2-\frac{L^2}{4}\right)\\
&=M\left(R^2-\frac{L^2}{6}\right).
\end{aligned}
\] 记 \[
K=R^2-\frac{L^2}{6},\qquad I_O=MK. \tag{2}
\]

\needspace{3\baselineskip}
\textbf{（二）小振动周期}

以最低位置为重力势能零点。杆偏转 \(\theta\) 后质心升高 \[
\Delta h=d-d\cos\theta=d(1-\cos\theta),
\] 故 \[
V=Mgd(1-\cos\theta).
\] 动能为 \[
\mathcal T=\frac12I_O\dot\theta^2.
\] 拉格朗日方程给出 \[
\frac{d}{dt}(I_O\dot\theta)+Mgd\sin\theta=0,
\] 即 \[
I_O\ddot\theta+Mgd\sin\theta=0. \tag{3}
\] 在 \(|\theta|\ll1\) 时，\(\sin\theta\simeq\theta\)： \[
\ddot\theta+\frac{Mgd}{I_O}\theta=0.
\] 因此 \[
\omega_0^2=\frac{gd}{K},
\] 周期为 \[
\boxed{
T=2\pi\sqrt{\frac{K}{gd}}
=2\pi\sqrt{
\frac{R^2-L^2/6}
{g\sqrt{R^2-L^2/4}}
}}. \tag{4}
\]

\begin{figure}[H]
\centering
\includegraphics[width=0.8\linewidth,height=\textheight,keepaspectratio,alt={图2-2　质心运动、重力和两端支持力}]{../figures/solution-02.pdf}
\caption*{图2-2　质心运动、重力和两端支持力}
\end{figure}

\needspace{3\baselineskip}
\textbf{（三）平衡点速度}

从转折点 \(\theta=\theta_m\) 静止释放，到平衡点 \(\theta=0\)
时，势能减少量全部转为动能： \[
\frac12I_O\dot\theta_0^2=Mgd(1-\cos\theta_m).
\] 所以 \[
\dot\theta_0^2
=\frac{2gd}{K}(1-\cos\theta_m)
\simeq\frac{gd}{K}\theta_m^2. \tag{5}
\]

\needspace{3\baselineskip}
\textbf{（四）用质心运动求碗的合力}

取 \(\bm e_r\) 沿 \(O\) 指向 \(G\)，\(\bm e_t\) 沿 \(\theta\)
增大方向。质心加速度为 \[
\bm a_G=-d\dot\theta^2\bm e_r+d\ddot\theta\bm e_t.
\] 重力在这组基矢中的分量是 \[
\bm W=Mg\cos\theta\,\bm e_r-Mg\sin\theta\,\bm e_t.
\] 设碗对杆的两个接触力的矢量和为 \(\bm N\)。由质心运动定理 \[
\bm N+\bm W=M\bm a_G,
\] 所以 \[
\begin{aligned}
N_r&=-Md\dot\theta^2-Mg\cos\theta,\\
N_t&=Md\ddot\theta+Mg\sin\theta.
\end{aligned} \tag{6}
\]

在平衡点 \(\theta=0\)，由对称性 \(N_t=0\)，合支持力大小为 \[
\begin{aligned}
N_0
&=Mg+Md\dot\theta_0^2\\
&\simeq Mg\left(1+\frac{d^2}{K}\theta_m^2\right).
\end{aligned} \tag{7}
\]

在转折点 \(\theta=\theta_m\)，有 \(\dot\theta=0\)。由式（3） \[
\ddot\theta_m=-\frac{gd}{K}\sin\theta_m.
\] 代入式（6）： \[
\begin{aligned}
N_r&=-Mg\cos\theta_m,\\
N_t&=Mg\sin\theta_m\left(1-\frac{d^2}{K}\right).
\end{aligned}
\] 因此合力大小 \[
N_m=Mg\sqrt{
\cos^2\theta_m+
\left(1-\frac{d^2}{K}\right)^2\sin^2\theta_m
}. \tag{8}
\] 令 \[
x=\frac{d^2}{K}.
\] 利用 \[
\cos^2\theta_m\simeq1-\theta_m^2,\qquad
\sin^2\theta_m\simeq\theta_m^2,
\] 式（8）的根号内为 \[
1+\big[(1-x)^2-1\big]\theta_m^2
=1+(-2x+x^2)\theta_m^2.
\] 再用 \(\sqrt{1+\varepsilon}\simeq1+\varepsilon/2\)： \[
N_m\simeq Mg\left[1+\left(-x+\frac{x^2}{2}\right)\theta_m^2\right]. \tag{9}
\] 由式（7）和式（9） \[
\begin{aligned}
\Delta N
&=N_m-N_0\\
&=Mg\left(-2x+\frac{x^2}{2}\right)\theta_m^2.
\end{aligned}
\] 故若``碗的支持力''指两个接触力的矢量合力大小，则 \[
\boxed{
\alpha=-2\frac{d^2}{K}
+\frac12\left(\frac{d^2}{K}\right)^2
},\qquad
d^2=R^2-\frac{L^2}{4},\quad
K=R^2-\frac{L^2}{6}. \tag{10}
\]

题面没有明确说明是否把``支持力''定义为两端法向力大小之和。若采用另一种约定
\(N=N_1+N_2\)，由两个法向力的径向分量 \[
\frac dR(N_1+N_2)=M(d\dot\theta^2+g\cos\theta)
\] 可得 \[
\boxed{
\alpha_{\rm sum}
=-\frac Rd\left(\frac12+\frac{d^2}{K}\right)
}. \tag{11}
\]
正式使用时应先说明所采用的定义；式（10）对应通常的``碗对杆的合力大小''。

\textbf{从牛顿方程直接求支持力合力。}
图2-2两端接触点均在圆周上，质心到圆心距离 \(d\) 恒定。取 \(\theta\) 为
\(OG\) 偏离竖直的角度，则 \[
\bm a_G=-d\dot\theta^2\bm e_r+d\ddot\theta\bm e_\theta.
\] 总有效惯量为 \(MK\)，\(K=d^2+L^2/12\)。角运动方程给 \[
\ddot\theta=-\frac{gd}{K}\sin\theta.
\] 重力的径向、切向分量为 \(Mg\cos\theta\) 与
\(-Mg\sin\theta\)。由牛顿第二定律，碗对杆合力的两分量是 \[
N_r=-M(d\dot\theta^2+g\cos\theta),
\] \[
N_\theta=Mg(1-d^2/K)\sin\theta.
\] 转折点 \(\theta=\theta_m\)
处速度为零，但切向加速度不为零。因此合力大小为 \[
N_m=Mg\sqrt{\cos^2\theta_m+(1-x)^2\sin^2\theta_m},\qquad x=d^2/K.
\] 在平衡点，能量守恒给 \(\dot\theta_0^2=2gd(1-\cos\theta_m)/K\)，故 \[
N_0=Mg[1+2x(1-\cos\theta_m)].
\] 展开至振幅二阶， \[
N_m=Mg[1+(-x+x^2/2)\theta_m^2]+O(\theta_m^4),
\] \[
N_0=Mg(1+x\theta_m^2)+O(\theta_m^4).
\] 相减得到正文系数
\(-2x+x^2/2\)。这个独立推导也说明为什么不能只看径向压力：转折点的切向支持力虽然为一阶，平方后正好贡献合力大小的二阶项。

若题意另取 \(N_1+N_2\)，必须将两接触力沿 \(OG\)
的投影相加后再比较，不能把上面的合力模长当成标量和。图中两个红箭头各指向圆心，其夹角不为零。

\textbf{一个可代回的杆长例子。} 取 \(R=2L\)，满足题设 \(R>1.5L\)。则 \[
d^2=R^2-L^2/4=15L^2/4,
\] \[
K=R^2-L^2/6=23L^2/6,
\] \[
x=d^2/K=45/46.
\] 支持力合力的二阶系数为 \[
\alpha=-2x+x^2/2=-\frac{6255}{4232}\simeq-1.4780.
\] 若改取两端力大小之和，系数为 \[
\alpha_{\rm sum}=-\frac{4}{\sqrt{15}}\left(\frac12+\frac{45}{46}\right)\simeq-1.5268.
\] 两者都是负值，但不同，说明仅检查正负并不足以排除支持力定义错误。

为核查展开精度，取 \(\theta_m=0.02\)。用精确合力公式分别计算 \(N_m\) 和
\(N_0\)，再除以
\(Mg\theta_m^2\)；随着振幅缩小一半，这个比值与二阶系数的误差应缩小约四倍，因为剩余项为
\(O(\theta_m^4)\)。

周期也可用无量纲数检查： \[
T\sqrt{g/L}=2\pi\sqrt{\frac{K/L^2}{d/L}}
=2\pi\sqrt{\frac{23/6}{\sqrt{15}/2}}.
\] 由于有效长度 \(K/d\) 具有长度单位，上式无量纲。若把 \(K\)
当成已含质量的转动惯量，却仍在公式中乘
\(M\)，量纲会多出一项质量，应在求数值之前修正。

\textbf{两端接触力各自的可实现性。}
合力正确之后，还要看两端是否都保留接触。取沿杆方向为
\(\bm e_\theta\)，两接触点相对圆心的矢量为
\(d\bm e_r\pm(L/2)\bm e_\theta\)。因此法向力沿各自指向圆心的方向，径向投影相同、切向投影相反：
\[
N_r=-\frac dR(N_1+N_2),
\qquad N_\theta=\frac L{2R}(N_2-N_1).
\] 解出两端支持力， \[
N_{1,2}=\frac R2\left(-\frac{N_r}{d}\mp\frac{2N_\theta}{L}\right).
\] 两者均非负等价于 \[
|N_\theta|\le\frac L{2d}|N_r|.
\] 这项条件不能仅从 \(N_1+N_2>0\)
推出。一个接触力变负意味着该端应脱离，原双接触运动方程失效。

在转折点，代入已求的两分量，得到 \[
(1-x)|\tan\theta_m|\le L/(2d),\qquad x=d^2/K.
\] 即 \[
|\tan\theta_m|\le\frac{6K}{dL}.
\] 原题的小振幅远小于该界限，因而保留双接触近似合理。对前面 \(R=2L\)
的例子，右侧约为
\(11.88\)；本例讨论的是小角展开，不能反过来把这个很大的界限当成二阶展开的适用角度。

平衡位置的速率不为零，两个接触力相等： \[
N_1=N_2=\frac{MgR}{2d}[1+2x(1-\cos\theta_m)].
\]
转折点一般不相等，但矢量合力仍须满足牛顿方程。图2-2的两个箭头应保留不同方向，不能用两个平行箭头来表示标量和。

对极短杆
\(L/R\to0\)，\(d\to R,K\to R^2,x\to1\)，质心运动趋于点质点在圆弧上的运动，切向支持力趋零。这时两端力之差的公式有
\(L\) 在分母，应先看 \(N_\theta\propto L^2\)
的抵消，不能直接宣称该差发散。

\newpage

\section{第 3 题}\label{ux7b2c-3-ux9898}

同心的内、外圆边界半径分别为 \(R_1\)、\(R_2\)，其中
\(R_1<R_2\)。环形区域中有垂直纸面向里的匀强磁场
\(B_0\)；内圆区域中另有垂直纸面的匀强磁场 \(B_x\)。规定垂直纸面向里为
\(B_x\) 的正方向。质量为 \(m\)、电荷量为 \(q>0\)
的粒子从外圆边界沿半径方向射入，初速度为 \(v_0\)。

（1）求粒子刚能触及内圆边界的临界初速度 \(v_{0\min}\)；

（2）若 \(v_0>v_{0\min}\)，调节 \(B_x\) 使粒子进入内圆后恰好通过圆心，求
\(B_x\)。

\begin{figure}[H]
\centering
\includegraphics[width=0.473\linewidth,height=\textheight,keepaspectratio,alt={图3-1　同心圆磁场区域与径向入射条件}]{../figures/question-03.pdf}
\caption*{图3-1　同心圆磁场区域与径向入射条件}
\end{figure}

\subsection{一、题型定位}\label{ux4e00ux9898ux578bux5b9aux4f4d-2}

临界入射轨迹与内圆相切。极坐标中的磁通项要在内外区域连续，才能把环区轨迹与中心穿越条件接起来。

\subsection{二、知识点}\label{ux4e8cux77e5ux8bc6ux70b9-2}

\begin{enumerate}
\def\labelenumi{\arabic{enumi}.}
\tightlist
\item
  洛伦兹力不做功，速率保持不变；
\item
  拉莫尔半径 \(\rho_c=mv/(qB)\)；
\item
  圆与同心圆相切的临界条件；
\item
  对称规范下的正则角动量；
\item
  分区磁场边界处机械速度连续；
\item
  有符号磁场方向的判定。
\end{enumerate}

\subsection{三、物理图像}\label{ux4e09ux7269ux7406ux56feux50cf-2}

正电荷沿半径向内进入指向纸内的磁场，洛伦兹力使轨迹向一侧弯曲。速度越大，轨迹圆半径越大，轨迹越``直''，越容易碰到内圆。临界情况下，环形区轨迹圆与内圆相切。第二问要求进入内圆后的圆弧穿过圆心，这意味着该圆弧相对于圆心的正则角动量为零。

\subsection{四、解题入口}\label{ux56dbux89e3ux9898ux5165ux53e3-2}

在环形磁场内速率始终为 \(v_0\)，轨迹曲率半径为 \[
\rho_c=\frac{mv_0}{qB_0}. \tag{1}
\]
入射点处速度沿半径，轨迹圆心在该点的切线方向，因此圆心到区域中心的距离是
\[
\sqrt{R_2^2+\rho_c^2}.
\] 轨迹到区域中心的最小距离便是该距离减去轨迹半径。

\subsection{五、解析}\label{ux4e94ux89e3ux6790-2}

\needspace{3\baselineskip}
\textbf{（一）进入内圆的临界几何}

轨迹到同心圆圆心 \(O\) 的最小距离为 \[
r_{\min}=\sqrt{R_2^2+\rho_c^2}-\rho_c. \tag{2}
\] 粒子能到达内圆边界要求 \(r_{\min}\le R_1\)。临界状态取等号： \[
\sqrt{R_2^2+\rho_c^2}-\rho_c=R_1.
\] 移项并平方： \[
\begin{aligned}
\sqrt{R_2^2+\rho_c^2}&=R_1+\rho_c,\\
R_2^2+\rho_c^2&=R_1^2+2R_1\rho_c+\rho_c^2,\\
\rho_c&=\frac{R_2^2-R_1^2}{2R_1}.
\end{aligned}
\] 结合式（1）得到 \[
\boxed{
v_{0\min}
=\frac{qB_0}{m}\frac{R_2^2-R_1^2}{2R_1}
}. \tag{3}
\]

\begin{figure}[H]
\centering
\includegraphics[width=0.8\linewidth,height=\textheight,keepaspectratio,alt={图3-2　环形区临界轨迹与正则角动量关系}]{../figures/solution-03.pdf}
\caption*{图3-2　环形区临界轨迹与正则角动量关系}
\end{figure}

\needspace{3\baselineskip}
\textbf{（二）环形区中的正则角动量}

取纸面向外为 \(z\)
正方向，而题目把向里的磁感应强度数值写成正数。因此环形区磁场矢量为 \[
\bm B=-B_0\bm e_z.
\] 选轴对称规范 \[
A_\varphi=-\frac12B_0r.
\] 角坐标 \(\varphi\) 是循环坐标，对应的正则角动量 \[
p_\varphi=mr^2\dot\varphi+qrA_\varphi
=L_z-\frac12qB_0r^2 \tag{4}
\] 在同一匀强磁场区域内守恒。

粒子刚进入外圆时速度沿半径，机械角动量为零，故 \[
p_\varphi=-\frac12qB_0R_2^2.
\] 粒子到达 \(r=R_1\) 的瞬间，其机械角动量为 \[
\begin{aligned}
L_z
&=p_\varphi+\frac12qB_0R_1^2\\
&=-\frac12qB_0(R_2^2-R_1^2).
\end{aligned} \tag{5}
\] 这个结果与 \(v_0\) 无关；\(v_0\) 只决定粒子能否到达该边界以及入射角。

\needspace{3\baselineskip}
\textbf{（三）内圆轨迹通过圆心的条件}

把内圆磁场 \(B_x\) 仍按``向里为正''处理。在内区 \[
p_\varphi^{(x)}=L_z-\frac12qB_xr^2.
\]
匀强磁场中的圆轨道若经过坐标原点，则相对于该原点的正则角动量为零。可由圆轨道几何直接验证：轨道经过原点时，任一点的切向机械角动量恰等于
\(qB_xr^2/2\)。所以在 \(r=R_1\) 处必须满足 \[
L_z-\frac12qB_xR_1^2=0. \tag{6}
\] 将式（5）代入： \[
-\frac12qB_0(R_2^2-R_1^2)
-\frac12qB_xR_1^2=0.
\] 约去 \(-q/2\) 后 \[
B_0(R_2^2-R_1^2)+B_xR_1^2=0.
\] 故 \[
\boxed{
B_x=-B_0\frac{R_2^2-R_1^2}{R_1^2}
}. \tag{7}
\] 负号表明内圆磁场实际应垂直纸面向外，与环形区磁场反向。

\textbf{严格进入与相切不同。} \(v_0=v_{0\min}\)
时只触及边界，没有进入内区的有限弧段；真正穿入需
\(v_0>v_{0\min}\)，临界值是这个开区间的下确界。

\textbf{全局规范也可保持连续。}
前面为简便分别使用两区的局部匀强场规范；这些表达式在环区还需加 \(C/r\)
项才能拼成一个全局矢势。取 \[
rA_\varphi=-\frac12\begin{cases}
B_xr^2,&r<R_1,\\
B_xR_1^2+B_0(r^2-R_1^2),&R_1<r<R_2.
\end{cases}
\] 此时同一个 \(p_\varphi=L_z+qrA_\varphi\)
在全程守恒并跨边界连续。两种算法都给式（7）；不能把``局部表达式的常数不同''说成正则角动量在轴对称静磁场中失守恒。

\textbf{用连续的磁通项检验中心穿越。} 取磁场正方向出纸面，环区场为
\(-B_0\)，内区场为 \(-B_x\)。旋转对称规范可取 \[
rA_\phi(r)=\int_0^r B_z(s)s\,ds.
\] 故内区、环区分别为 \[
rA_\phi=-\frac12B_xr^2,
\] \[
rA_\phi=-\frac12[B_xR_1^2+B_0(r^2-R_1^2)].
\] 两式在 \(R_1\) 连续。正则角动量 \[
p_\phi=mr^2\dot\phi+qrA_\phi
\]
在整个轴对称区域守恒，不能在每个区域另取不连续的零势基准后，把两个不同常数拼接起来。

要穿过 \(r=0\)，有限速度轨迹必须有
\(p_\phi=0\)。初始径向入射时机械角动量为零，所以外边界处应满足 \[
B_xR_1^2+B_0(R_2^2-R_1^2)=0.
\] 因此 \[
B_x=-B_0\frac{R_2^2-R_1^2}{R_1^2}.
\] 这个符号要求中心磁场与环区反向；题面图中的 \(\otimes B_x\)
只是正参数方向标记，解出负 \(B_x\) 后实际应为出纸面。

当 \(p_\phi=0\) 时，机械切向速度由 \(-qA_\phi/m\)
决定。环区的最大切向速度出现在内边界，内区则随 \(r\)
减小而减小。只要入射速率严格超过环区切向速度的最大值，径向动能始终为正，粒子可继续到中心。临界相切情况只接触边界，不能用等号代替进入内圆的条件。

\textbf{给定入射速率后的场强例子。} 取 \(R_2=2R_1\)。临界回旋半径为 \[
\rho_{\rm crit}=\frac{4R_1^2-R_1^2}{2R_1}=\frac32R_1.
\] 临界速率为 \(|q|B_0\rho_{\rm crit}/m\)。若选
\(v_0=2|q|B_0R_1/m\)，则严格超过阈值；进入内圆后要过中心，中心场必须取
\[
B_x=-3B_0.
\] 全局磁通项在外边界为零，故正则角动量为零。环区切向速率为 \[
v_\phi(r)=\frac{qB_0}{2m}\left(r-\frac{4R_1^2}{r}\right),
\] 符号取决于电荷；其绝对值在内边界为
\(3|q|B_0R_1/(2m)\)。径向速率由总速率守恒得到 \[
v_r^2=v_0^2-v_\phi^2.
\] 在内边界，\(v_r^2=(7/4)(qB_0R_1/m)^2>0\)，所以不是擦边。

内区的切向速率为 \(v_\phi=qB_xr/(2m)\)，其大小随 \(r\)
下降；因此已进入内区后没有新的离心势垒，最终可以穿过
\(r=0\)。在中心，切向速率为零，速度方向沿瞬时径向通过中心。

场强 \(B_x\) 与 \(v_0\)
无关不是漏项：中心穿越所需的角动量条件由总磁通决定，而速度只负责保证有足够径向动能越过环区。这两类条件分别对应轨道方向与可达性。

\textbf{内圆的通过时间与径向动能。}
中心穿越所需场强已由磁通确定，进入内圆后的时间则还依赖速率。取
\(\omega_c=qB_x/m\)，其符号由场和电荷决定。在 \(p_\phi=0\)
的连续规范下，内区切向速率为 \[
v_\phi=\omega_c r/2.
\] 磁力不做功，故 \[
v_r^2=v_0^2-\omega_c^2r^2/4.
\] 由内边界走到中心的时间为 \[
t_{1\to0}=\int_0^{R_1}\frac{dr}{\sqrt{v_0^2-\omega_c^2r^2/4}}
=\frac2{|\omega_c|}\arcsin\frac{|\omega_c|R_1}{2v_0}.
\] 只有参数在允许区内，反正弦才有意义。前面
\(R_2=2R_1,v_0=2|q|B_0R_1/m\) 的例子给 \[
\frac{|\omega_c|R_1}{2v_0}=\frac34,
\qquad t_{1\to0}=\frac{2m}{3|q|B_0}\arcsin(3/4).
\] 在这一段中 \(\dot\phi=v_\phi/r=\omega_c/2\) 为常量，但 \(\dot r\)
不恒定。极角相对入口的变化为
\(\omega_ct/2\)；过中心时极坐标角本身出现坐标跳变，真实速度和轨道保持连续。

把速度提高至很大，反正弦展开给
\(t_{1\to0}\to R_1/v_0\)，恢复近似直线飞行时间。若在阈值附近，内边界径向速度很小，但内区积分的平方根端点奇性可积，不能只凭
\(v_r(R_1)=0\) 就认定后续时间无限。

严格阈值仍由环区的实际进入几何判断。精确临界圆轨道在边界相切并返回，没有进入内区；上面的积分只描述已经穿入内区的支路。阈值以上的轨迹可接近积分的边界值，这与``等号轨迹不能进入''不矛盾。

图3-2应同时核查三件事：入射速率保证到达内边界，内外总磁通保证
\(p_\phi=0\)，中心实际磁场方向反转。只调整其中一项不能保证中心穿越。

\newpage

\section{第 4 题}\label{ux7b2c-4-ux9898}

（1）半径为 \(r\)、质量为 \(m\) 的均匀圆环在水平地面上以质心速度 \(v_1\)
向右纯滚动，撞到高度为 \(h\) 的台阶，其中
\(h<r/2\)。碰撞完全非弹性，求圆环能够翻过台阶并继续向前滚动时 \(v_1\)
应满足的条件。

（2）可变形圆环纯滚动时受到滚动阻力偶矩 \(M=kN\)，其中 \(N\)
为地面对圆环的压力，\(k\) 为已知长度。若初速度为 \(v_0\)，求停止时间。

（3）在摩擦足够大的水平地面上，一根铁棍从左侧推动上述圆环，使其向右做匀速纯滚动。铁棍与圆环间的动摩擦因数为
\(\mu\)，忽略二者间的滚动摩擦。设接触点半径与竖直向下方向夹角为
\(\theta\)。求使铁棍推力最小时的 \(\theta\) 及最小推力。

\begin{figure}[H]
\centering
\includegraphics[width=0.8\linewidth,height=\textheight,keepaspectratio,alt={图4-1　滚环与台阶碰撞的几何关系}]{../figures/question-04.pdf}
\caption*{图4-1　滚环与台阶碰撞的几何关系}
\end{figure}

\subsection{一、题型定位}\label{ux4e00ux9898ux578bux5b9aux4f4d-3}

越台阶先按接触角点写碰撞冲量，再检查碰后的机械能。恒速推环则需同时平衡力与力矩，并保留地面支持力为正的条件。

\subsection{二、知识点}\label{ux4e8cux77e5ux8bc6ux70b9-3}

\begin{enumerate}
\def\labelenumi{\arabic{enumi}.}
\tightlist
\item
  瞬时碰撞中关于冲量作用点的角动量守恒；
\item
  圆环转动惯量 \(I_C=mr^2\)；
\item
  越障的临界能量条件；
\item
  滚动阻力矩的功率；
\item
  刚体平动、转动平衡；
\item
  三角线性组合的最大值。
\end{enumerate}

\subsection{三、物理图像}\label{ux4e09ux7269ux7406ux56feux50cf-3}

撞上台阶棱角 \(A\) 后，原接地点失去约束，圆环瞬时改为绕 \(A\)
转动。碰撞过程中 \(A\) 点冲量很大，但它对 \(A\)
的力矩为零；重力冲量可忽略，所以只能关于 \(A\)
守恒角动量，不能守恒机械能。越障阶段则由碰撞后的转动动能支付质心上升
\(h\) 所需的重力势能。

\subsection{四、解题入口}\label{ux56dbux89e3ux9898ux5165ux53e3-3}

碰撞前圆环既有质心平动角动量，也有绕质心自转角动量。相对于
\(A\)，两项方向相同，大小分别是 \[
mv_1(r-h),\qquad I_C\frac{v_1}{r}=mrv_1.
\] 碰撞后圆环绕 \(A\) 的转动惯量为 \[
I_A=I_C+mr^2=2mr^2.
\]

\subsection{五、解析}\label{ux4e94ux89e3ux6790-3}

\needspace{3\baselineskip}
\textbf{（一）碰撞后的角速度}

设碰撞后绕 \(A\) 的角速度大小为 \(\Omega\)。关于 \(A\) 写角动量守恒： \[
mv_1(r-h)+mrv_1=I_A\Omega.
\] 代入 \(I_A=2mr^2\)： \[
\begin{aligned}
mv_1(2r-h)&=2mr^2\Omega,\\
\Omega&=\frac{v_1(2r-h)}{2r^2}.
\end{aligned} \tag{1}
\] 碰撞后的转动动能为 \[
\begin{aligned}
E_+
&=\frac12I_A\Omega^2\\
&=\frac12(2mr^2)
\left[\frac{v_1(2r-h)}{2r^2}\right]^2\\
&=\frac{mv_1^2(2r-h)^2}{4r^2}.
\end{aligned} \tag{2}
\] 质心从碰撞位置转到 \(A\) 的正上方时升高 \(h\)。临界越障条件为 \[
E_+=mgh.
\] 代入式（2）并约去 \(m\)： \[
\frac{v_1^2(2r-h)^2}{4r^2}=gh.
\] 所以 \[
\boxed{
v_1>\frac{2r\sqrt{gh}}{2r-h}
}. \tag{3}
\]

等号只给能量阈值：圆心在棱角正上方处速度趋零，该点是势垒顶，理想运动渐近接近它而不会在有限时间翻过去。题目要求翻过后继续滚动，故需严格大于阈值。这个式子是能量下限；下文``保持角点接触''一节进一步给出同一运动支路的速度上限，合并后才得到完整的接触越障区间。

\needspace{3\baselineskip}
\textbf{（二）滚动阻力导致的停止时间}

圆环纯滚动时 \(\omega=v/r\)。总动能为 \[
\begin{aligned}
E
&=\frac12mv^2+\frac12I_C\omega^2\\
&=\frac12mv^2+\frac12mr^2\frac{v^2}{r^2}\\
&=mv^2.
\end{aligned} \tag{4}
\] 水平面上 \(N=mg\)，滚动阻力偶矩大小为 \[
M=kmg.
\] 其耗散功率为 \[
\frac{dE}{dt}=-M\omega=-kmg\frac vr.
\] 又 \[
\frac{dE}{dt}=\frac{d(mv^2)}{dt}=2mv\frac{dv}{dt}.
\] 当 \(v>0\) 时约去 \(mv\)： \[
2\frac{dv}{dt}=-\frac{kg}{r},
\] 即 \[
\frac{dv}{dt}=-\frac{kg}{2r}. \tag{5}
\] 这是恒定减速度。由 \(0=v_0-(kg/2r)t_{\rm stop}\) 得 \[
\boxed{
t_{\rm stop}=\frac{2rv_0}{kg}
}. \tag{6}
\]

\needspace{3\baselineskip}
\textbf{（三）推环的受力与力矩平衡}

如图4-2，铁棍作用在圆环左下侧。设铁棍对圆环的法向力大小为
\(P_n\)，接触处滑动摩擦力大小为
\(\mu P_n\)。按图示方向分解，铁棍对圆环的水平、竖直分量为 \[
\begin{aligned}
P_x&=P_n(\sin\theta+\mu\cos\theta),\\
P_y&=P_n(\cos\theta-\mu\sin\theta).
\end{aligned} \tag{7}
\] 推力合力大小为 \[
P=P_n\sqrt{1+\mu^2}. \tag{8}
\]

\begin{figure}[H]
\centering
\includegraphics[width=0.82\linewidth,height=\textheight,keepaspectratio,alt={图4-2　铁棍推动匀速滚环时的受力与力矩}]{../figures/solution-04.pdf}
\caption*{图4-2　铁棍推动匀速滚环时的受力与力矩}
\end{figure}

设地面摩擦力向右为正，记为 \(f\)，地面支持力为 \(N\)。匀速平动给出 \[
\begin{aligned}
P_x+f&=0,\\
N+P_y-mg&=0.
\end{aligned}
\] 所以 \[
\begin{aligned}
f&=-P_n(\sin\theta+\mu\cos\theta),\\
N&=mg-P_n(\cos\theta-\mu\sin\theta).
\end{aligned} \tag{9}
\] 取逆时针力矩为正。地面摩擦、棍环摩擦以及滚动阻力偶矩共同满足 \[
rf+r\mu P_n+kN=0. \tag{10}
\] 将式（9）代入： \[
\begin{aligned}
0={}&-rP_n(\sin\theta+\mu\cos\theta)
+r\mu P_n\\
&+k\big[mg-P_n(\cos\theta-\mu\sin\theta)\big].
\end{aligned}
\] 把含 \(P_n\) 的各项移到同一侧： \[
P_n\big[
(r-k\mu)\sin\theta
+(r\mu+k)\cos\theta-r\mu
\big]=kmg. \tag{11}
\] 因此 \[
P=\frac{kmg\sqrt{1+\mu^2}}
{(r-k\mu)\sin\theta+(r\mu+k)\cos\theta-r\mu}. \tag{12}
\] 要使 \(P\) 最小，应使分母最大。令 \[
A=r-k\mu,\qquad B=r\mu+k.
\] 当 \(A>0\) 且最优角位于允许区间内部时， \[
A\sin\theta+B\cos\theta
\le\sqrt{A^2+B^2},
\] 等号条件为 \[
\boxed{
\tan\theta_{\rm opt}=\frac{r-k\mu}{r\mu+k}
}. \tag{13}
\] 注意 \[
\begin{aligned}
A^2+B^2
&=(r-k\mu)^2+(r\mu+k)^2\\
&=(1+\mu^2)(r^2+k^2).
\end{aligned}
\] 故 \[
\boxed{
P_{\min}
=\frac{kmg\sqrt{1+\mu^2}}
{\sqrt{(1+\mu^2)(r^2+k^2)}-r\mu}
}. \tag{14}
\]
若参数使式（13）给出的角不在实际可接触的角度范围内，则最小值应在该范围的端点比较，不能机械套用内部极值。

\textbf{碰撞耗能与越障阈值的分段检查。} 圆环碰前绕自身圆心的惯量为
\(mr^2\)，其纯滚动总动能为 \(mv_0^2\)。碰撞瞬间对角点 \(A\)
取矩，碰前角动量大小为 \[
L_A^-=mv_0(r-h)+mr^2(v_0/r)=mv_0(2r-h).
\] 碰后绕 \(A\) 的惯量是 \(I_A=2mr^2\)，因此 \[
\Omega_+=\frac{v_0(2r-h)}{2r^2}.
\] 碰后动能为 \[
K_+=\frac12I_A\Omega_+^2
=\frac{mv_0^2(2r-h)^2}{4r^2}.
\]
与碰前相比，机械能减少；碰撞期间不能直接使用能量守恒。之后上摆若保持接触且无进一步耗散，圆心最高时需增加势能
\(mgh\)，因而有限时间越过要求 \[
K_+>mgh.
\]
它恢复正文的严格速度阈值。还应由径向牛顿方程逐角核查接触力非负；过大的速度可能使环与角点脱离，需改用抛体运动。

平地滚动阻力的检查则不涉及撞击。由 \(E=mv^2\)、阻力功率 \(-kmgv/r\)，有
\[
2mv\dot v=-\frac{kmg}{r}v.
\] 在 \(v>0\) 区间约去 \(v\) 后，减速度恒为 \(kg/(2r)\)，停止时间 \[
t_{\rm stop}=\frac{2rv_0}{kg}.
\] 用恒定减速度的路程公式还可得
\(s_{\rm stop}=rv_0^2/(kg)\)。两个结果同时检查了滚动动能的系数二。推环问题则在图4-2上另列棍的法向、切向力和地面力，不能把滚动阻力偶矩当成一股沿地面作用的力。

\textbf{圆盘越台阶的完整对照。}
将薄环换为均匀圆盘，\(I_C=mr^2/2\)。碰前角点角动量为 \[
L_A^-=mv_0(r-h)+\frac12mr^2\frac{v_0}{r}
=mv_0(3r/2-h).
\] 碰后惯量为 \(I_A=3mr^2/2\)，故 \[
\Omega_+=\frac{v_0(3r/2-h)}{3r^2/2}.
\] 碰后动能为 \[
K_+=\frac{mv_0^2(3r/2-h)^2}{3r^2}.
\] 要求 \(K_+>mgh\)，得到 \[
v_0>\frac{\sqrt3\,r\sqrt{gh}}{3r/2-h}.
\] 这与薄环的阈值不同，因为自转动量和碰后惯量都改变了。取
\(h=r/2\)，圆盘阈值为 \(\sqrt{3gr/2}\)，薄环阈值为
\(\sqrt{8gr/9}\)；同一台阶下薄环较易保留绕角点的转动能。

平地滚动阻力偶矩若仍为 \(kmg\)，圆盘总能为 \(3mv^2/4\)，因此 \[
\dot v=-\frac{2kg}{3r},\qquad
 t_{\rm stop}=\frac{3rv_0}{2kg}.
\] 对任意 \(I_C=\kappa mr^2\)，两种问题可统一写为 \[
I_A=(1+\kappa)mr^2,
\qquad L_A^-=mv_0[(1+\kappa)r-h],
\] \[
v_{\rm th}=\frac{r\sqrt{2(1+\kappa)gh}}{(1+\kappa)r-h}.
\] 代入 \(\kappa=1\)、\(1/2\)
应分别恢复薄环和圆盘公式；分母非正时，碰后旋转方向也需重判，不能直接沿用越障支路。

\textbf{保持角点接触给出的速度上限。}
能量下限还不是完整的接触运动条件。令圆心到角点的连线与竖直向上方向夹角为
\(\psi\)，初始位置有 \[
\cos\psi_0=1-h/r.
\] 碰后能量守恒给 \[
r^2\dot\psi^2=K_+/m-gr(\cos\psi-\cos\psi_0).
\] 沿 \(AC\) 外法向写平动方程，角点支持力为 \[
N=mg\cos\psi-mr\dot\psi^2
=mg(2\cos\psi-\cos\psi_0)-K_+/r.
\] 上摆时 \(\cos\psi\) 增大，所以最小支持力在刚碰后。保持接触要求 \[
N_0\ge0\quad\Longleftrightarrow\quad K_+\le mg(r-h).
\] 与有限时间越障的下限合并，得到在角点约束可实现的模型中 \[
mgh<K_+\le mg(r-h),
\] \[
\frac{2r\sqrt{gh}}{2r-h}<v_1
\le\frac{2r\sqrt{g(r-h)}}{2r-h}.
\] 上下界之间存在区间恰好要求
\(h<r/2\)，与题设一致。若只算下限，就看不出这个高度限制的接触含义。

下限等号是势垒顶的渐近轨迹，必须排除；上限等号刚碰后 \(N=0\)，之后上摆
\(N\)
增大，可作为无滑动接触约束的边界。维持绕角点纯滚动还需要足够的切向约束，题目的完全非弹性接触模型默认这一点。

速度超过上限后，支持力公式要求负值，实际圆环应离开角点。之后可能以抛体运动越过台阶，但不能沿用绕角点的轨迹和能量阈值。若要讨论该支路，须另求离开时速度、与台阶的再次接触及其碰撞规则；本文的区间明确对应保持角点接触的越障支路。

\newpage

\section{第 5 题}\label{ux7b2c-5-ux9898}

三根竖直试管的横截面积依次为
\(S\)、\(2S\)、\(S\)，底部连通且连通部分体积可忽略。管中装有密度为
\(\rho\) 的理想流体，平衡时液柱高度均为
\(l\)。忽略黏滞，并假定同一根试管内各处流速相同。小扰动后三根液柱高度分别为
\[
l[1+x_1(t)],\qquad l[1+x_2(t)],\qquad l[1+x_3(t)],
\] 其中 \(|x_i|\ll1\)。

（1）写出 \(x_1,x_2,x_3\) 的约束关系；

（2）用 \(x_1,x_3,\dot x_1,\dot x_3\) 表示系统的动能和势能；

（3）初态为 \(x_1=x_2=\varepsilon\)、系统静止，其中
\(\varepsilon\ll1\)，求 \(x_1(t)\)。

\begin{figure}[H]
\centering
\includegraphics[width=0.8\linewidth,height=\textheight,keepaspectratio,alt={图5-1　三根不同截面积连通管的液面位移}]{../figures/question-05.pdf}
\caption*{图5-1　三根不同截面积连通管的液面位移}
\end{figure}

\subsection{一、题型定位}\label{ux4e00ux9898ux578bux5b9aux4f4d-4}

三个液面位移只有两个独立量。体积守恒与流体动能共同给出小振动方程，本题的两个模态频率相同。

\subsection{二、知识点}\label{ux4e8cux77e5ux8bc6ux70b9-4}

\begin{enumerate}
\def\labelenumi{\arabic{enumi}.}
\tightlist
\item
  不可压缩流体的体积守恒；
\item
  均匀流速近似下液柱动能；
\item
  重力势能的二阶展开；
\item
  拉格朗日方程；
\item
  质量矩阵与刚度矩阵；
\item
  简谐振动初值问题。
\end{enumerate}

\subsection{三、物理图像}\label{ux4e09ux7269ux7406ux56feux50cf-4}

左、中、右三管液面分别升高
\(lx_1,lx_2,lx_3\)。中管面积是两侧的两倍，所以同样的无量纲位移会带来两倍体积变化。流体不可压缩，三个体积变化的代数和必须为零。每根管内液体可近似成整体以速度
\(l\dot x_i\) 上下运动。

\subsection{四、解题入口}\label{ux56dbux89e3ux9898ux5165ux53e3-4}

体积守恒立即给出 \[
Slx_1+2Slx_2+Slx_3=0,
\] 即 \[
\boxed{x_1+2x_2+x_3=0},\qquad
x_2=-\frac{x_1+x_3}{2}. \tag{1}
\] 用式（1）消去 \(x_2\) 后，系统只剩 \(x_1,x_3\) 两个广义坐标。

\subsection{五、解析}\label{ux4e94ux89e3ux6790-4}

\needspace{3\baselineskip}
\textbf{（一）动能}

在二阶小振动近似中，各液柱质量可取平衡值： \[
m_1=\rho Sl,\qquad m_2=2\rho Sl,\qquad m_3=\rho Sl.
\] 三根液柱的速度分别为 \[
v_i=l\dot x_i.
\] 所以总动能是 \[
\begin{aligned}
\mathcal T
&=\frac12\rho Sl(l\dot x_1)^2
+\frac12(2\rho Sl)(l\dot x_2)^2
+\frac12\rho Sl(l\dot x_3)^2\\
&=\frac12\rho Sl^3
\left(\dot x_1^2+2\dot x_2^2+\dot x_3^2\right).
\end{aligned} \tag{2}
\] 由式（1） \[
\dot x_2=-\frac{\dot x_1+\dot x_3}{2}.
\] 逐项展开： \[
\begin{aligned}
\dot x_1^2+2\dot x_2^2+\dot x_3^2
&=\dot x_1^2+\frac12(\dot x_1+\dot x_3)^2+\dot x_3^2\\
&=\dot x_1^2+\frac12\dot x_1^2
+\dot x_1\dot x_3+\frac12\dot x_3^2+\dot x_3^2\\
&=\frac12\left(3\dot x_1^2+2\dot x_1\dot x_3+3\dot x_3^2\right).
\end{aligned}
\] 因此 \[
\boxed{
\mathcal T=\frac{\rho Sl^3}{4}
\left(3\dot x_1^2+2\dot x_1\dot x_3+3\dot x_3^2\right)
}. \tag{3}
\]

\needspace{3\baselineskip}
\textbf{（二）重力势能}

截面积为 \(A\)、高度为 \(h\) 的液柱重力势能相对于管底为 \[
V_h=(\rho Ah)g\frac h2=\frac12\rho gAh^2.
\] 因此 \[
V=\frac12\rho gSl^2
\left[(1+x_1)^2+2(1+x_2)^2+(1+x_3)^2\right].
\] 展开括号。常数项不影响运动；一次项为 \[
2(x_1+2x_2+x_3)=0
\] ，恰由体积约束消失。于是只保留二次项： \[
V=\frac12\rho gSl^2
\left(x_1^2+2x_2^2+x_3^2\right)+{\rm const}.
\] 用式（1）消去 \(x_2\)： \[
\boxed{
V=\frac{\rho gSl^2}{4}
\left(3x_1^2+2x_1x_3+3x_3^2\right)+{\rm const}
}. \tag{4}
\]

\needspace{3\baselineskip}
\textbf{（三）运动方程}

对 \(x_1\) 使用拉格朗日方程。由式（3） \[
\frac{\partial\mathcal T}{\partial\dot x_1}
=\frac{\rho Sl^3}{2}(3\dot x_1+\dot x_3),
\] 所以 \[
\frac{d}{dt}\frac{\partial\mathcal T}{\partial\dot x_1}
=\frac{\rho Sl^3}{2}(3\ddot x_1+\ddot x_3).
\] 由式（4） \[
\frac{\partial V}{\partial x_1}
=\frac{\rho gSl^2}{2}(3x_1+x_3).
\] 从而 \[
3\ddot x_1+\ddot x_3+\frac gl(3x_1+x_3)=0. \tag{5}
\] 同理对 \(x_3\)： \[
\ddot x_1+3\ddot x_3+\frac gl(x_1+3x_3)=0. \tag{6}
\] 把式（5）、（6）写成矩阵形式： \[
\begin{pmatrix}3&1\\1&3\end{pmatrix}
\begin{pmatrix}\ddot x_1\\\ddot x_3\end{pmatrix}
+\frac gl
\begin{pmatrix}3&1\\1&3\end{pmatrix}
\begin{pmatrix}x_1\\x_3\end{pmatrix}=0.
\] 矩阵可逆，两边左乘其逆矩阵，得到 \[
\boxed{
\ddot x_1+\frac glx_1=0,\qquad
\ddot x_3+\frac glx_3=0
}. \tag{7}
\]
这并不是说两根液柱彼此没有耦合，而是动能矩阵与势能矩阵具有相同形状，约化后所有允许运动都具有同一个频率
\[
\omega=\sqrt{\frac gl}. \tag{8}
\]

\begin{figure}[H]
\centering
\includegraphics[width=0.536\linewidth,height=\textheight,keepaspectratio,alt={图5-2　给定初态下三液面的同频运动}]{../figures/solution-05.pdf}
\caption*{图5-2　给定初态下三液面的同频运动}
\end{figure}

\needspace{3\baselineskip}
\textbf{（四）代入初始条件}

初态 \(x_1(0)=x_2(0)=\varepsilon\)。由体积约束 \[
\varepsilon+2\varepsilon+x_3(0)=0,
\] 所以 \[
x_3(0)=-3\varepsilon.
\] 系统静止意味着 \[
\dot x_1(0)=\dot x_2(0)=\dot x_3(0)=0.
\] 式（7）的第一个方程通解为 \[
x_1(t)=A\cos\omega t+B\sin\omega t.
\] 由 \(x_1(0)=\varepsilon\) 得 \(A=\varepsilon\)；由 \(\dot x_1(0)=0\)
得 \(B=0\)。故 \[
\boxed{
x_1(t)=\varepsilon\cos\left(\sqrt{\frac gl}\,t\right)
}. \tag{9}
\] 同时 \[
x_2(t)=\varepsilon\cos\left(\sqrt{\frac gl}\,t\right),\qquad
x_3(t)=-3\varepsilon\cos\left(\sqrt{\frac gl}\,t\right),
\] 可直接检验体积约束在任意时刻均成立。

\textbf{用约束矩阵复查两种模态。} 设
\(\bm q=(x_1,x_3)^T\)。消去中管位移后，动能和势能分别为 \[
\mathcal T=\frac12\dot{\bm q}^{T}\bm M\dot{\bm q},\qquad
V=\frac12\bm q^T\bm K\bm q+V_0,
\] 其中 \[
\bm M=\frac{\rho Sl^3}{2}\begin{pmatrix}3&1\\1&3\end{pmatrix},
\qquad
\bm K=\frac{\rho gSl^2}{2}\begin{pmatrix}3&1\\1&3\end{pmatrix}.
\] 两矩阵成比例，所以 \[
\bm M^{-1}\bm K=(g/l)\bm I.
\] 任何满足体积约束的位移方向都以同一频率振动。可选对称模态
\((x_1,x_2,x_3)=(1,-1,1)\) 与反对称模态 \((1,0,-1)\)，其频率均为
\(\sqrt{g/l}\)。

题给初态为 \((1,1,-3)\varepsilon\)，可分解为 \[
(1,1,-3)=-(1,-1,1)+2(1,0,-1).
\] 两模态同时乘
\(\cos\omega t\)，便得到图5-2的振幅比。初始静止意味着所有正弦项系数为零。

势能的一次项应恰好消失： \[
\rho gSl^2(x_1+2x_2+x_3)=0.
\]
若消去变量后仍留下非零一次项，说明参考液面或体积约束写错。动能矩阵的本征值与
\(4,2\) 成比例，均为正，也排除了伪自由度。

频率简并依赖``每根管的平衡高度相同、管内速度均匀、连接段体积和惯量忽略''三项条件。若连接段较长，其动能会增加质量矩阵但不同比例增加势能矩阵，两频率一般分裂。不能把这里的简并推广为任意三根连通管的性质。

\textbf{连接段惯量破坏简并的带解例子。}
仍取原三管，但设连接段只对反对称流动增加动能
\(\frac12m_c l^2(\dot x_1-\dot x_3)^2\)，以这一指定简化模型说明频率变化。定义对称、反对称坐标
\[
u=(x_1+x_3)/\sqrt2,\qquad w=(x_1-x_3)/\sqrt2.
\] 原动能为 \[
\mathcal T=\frac{\rho Sl^3}{2}(2\dot u^2+\dot w^2),
\] 原势能为 \(\rho gSl^2(2u^2+w^2)/2\)。连接段增加
\(m_cl^2\dot w^2\)，因此 \[
\omega_u^2=g/l,
\qquad
\omega_w^2=\frac{g/l}{1+2m_c/(\rho Sl)}.
\] 对称频率不变，反对称频率变小。取 \(m_c=\rho Sl/2\)，则
\(\omega_w=\omega_u/\sqrt2\)。

原初态给
\(u(0)=-\sqrt2\varepsilon\)、\(w(0)=2\sqrt2\varepsilon\)、两初速度为零。于是
\[
u(t)=-\sqrt2\varepsilon\cos\omega_ut,
\qquad w(t)=2\sqrt2\varepsilon\cos\omega_wt.
\] 换回液面， \[
x_1(t)=-\varepsilon\cos\omega_ut+2\varepsilon\cos\omega_wt,
\] \[
x_3(t)=-\varepsilon\cos\omega_ut-2\varepsilon\cos\omega_wt,
\qquad x_2(t)=\varepsilon\cos\omega_ut.
\] 逐时仍有
\(x_1+2x_2+x_3=0\)。这是一项明确指定额外动能的变式，不是声称实际任意形状连接段都只影响一种模态；真实连接段需先由流量分布计算附加质量。

\textbf{初态能量如何在三根管之间分配。} 原初态为
\((x_1,x_2,x_3)=\varepsilon(1,1,-3)\)，速度为零。位移均为无量纲量，实际液面变化是
\(lx_i\)。忽略连接段惯量时，总振动能为 \[
E=\frac12\rho gSl^2(x_1^2+2x_2^2+x_3^2)
=6\rho gSl^2\varepsilon^2.
\] 运动中三位移同乘 \(\cos\omega t\)，所以 \[
V(t)-V_0=E\cos^2\omega t,
\qquad\mathcal T(t)=E\sin^2\omega t.
\] 最低势能时三液面同时通过平衡位置，速度比例仍为
\(1:1:-3\)。体积流量比例为 \[
S\dot h_1:2S\dot h_2:S\dot h_3=1:2:-3,
\] 三者和为零。这一流量检查与液面位移比例不同，不能遗漏中管的截面积二。

在模态坐标 \(u=(x_1+x_3)/\sqrt2,w=(x_1-x_3)/\sqrt2\) 中，初值是
\(u_0=-\sqrt2\varepsilon,w_0=2\sqrt2\varepsilon\)。两独立模态能量为 \[
E_u=\rho gSl^2u_0^2=2\rho gSl^2\varepsilon^2,
\] \[
E_w=\frac12\rho gSl^2w_0^2=4\rho gSl^2\varepsilon^2.
\]
相加恢复总能。反对称模态的坐标幅值较大，但不能直接按幅值平方比四比一分配能量，因为两坐标的有效惯量不同。

如果只改变初速度，通解仍为每一坐标的正弦与余弦组合： \[
x_i(t)=x_i(0)\cos\omega t+\frac{\dot x_i(0)}\omega\sin\omega t.
\]
初位移和初速度都必须满足体积约束。图5-2的三条同相曲线仅对应原题静止释放初态；任意初速度可能使三液面相位不同，但不会在该理想模型中分裂频率。

\newpage

\section{第 6 题}\label{ux7b2c-6-ux9898}

温度为 \(T\) 的平衡态理想气体满足麦克斯韦速度分布 \[
dP=\left(\frac{m}{2\pi kT}\right)^{3/2}
\exp\left(-\frac{m v^2}{2kT}\right)
dv_x\,dv_y\,dv_z.
\] 一根水平圆柱形试管长为 \(l\)、横截面积为 \(S\)、质量为
\(M\)，其中密封着分子质量为 \(m\)、温度为 \(T\)、压强为 \(p_0\)
的理想气体。某时刻突然去掉左壁，保留右壁。定义 \[
f(x)=\int_0^x e^{-u^2}\,du.
\]

（1）试管固定，求时刻 \(t\) 前累计与右壁碰撞的分子数；

（2）试管固定，右壁碰撞完全弹性，求时刻 \(t\) 前分子传给右壁的总动量；

（3）试管可自由运动且很重，求最终速度；

（4）把气体换为与温度为 \(T\)
的试管保持热平衡的光子气体，光子与右壁完全弹性碰撞。已知黑体辐射出射度
\(J=\sigma T^4\)，求很重试管的最终速度。

\begin{figure}[H]
\centering
\includegraphics[width=0.8\linewidth,height=\textheight,keepaspectratio,alt={图6-1　左端开启、右壁保留的气体试管}]{../figures/question-06.pdf}
\caption*{图6-1　左端开启、右壁保留的气体试管}
\end{figure}

\subsection{一、题型定位}\label{ux4e00ux9898ux578bux5b9aux4f4d-5}

原题统计的是撞过右壁的分子数，只有初始向右分子参与。按初始位置与正速度积分后，再增加单次冲量权重求总动量；这与全部分子的逸出计数不同。

\subsection{二、知识点}\label{ux4e8cux77e5ux8bc6ux70b9-5}

\begin{enumerate}
\def\labelenumi{\arabic{enumi}.}
\tightlist
\item
  麦克斯韦分布的一维边缘分布；
\item
  初始位置均匀分布；
\item
  分段碰撞条件；
\item
  完全弹性反射的冲量；
\item
  动量守恒；
\item
  黑体辐射能量密度 \(u=4\sigma T^4/c\)；
\item
  各向同性分布的 \(\langle|\cos\theta|\rangle=1/2\)。
\end{enumerate}

\subsection{三、物理图像}\label{ux4e09ux7269ux7406ux56feux50cf-5}

左壁撤去后，\(v_x<0\) 的分子直接从左端逃逸；\(v_x>0\)
的分子若来得及先撞右壁，就反射后再从左端逃逸。对给定正速度
\(v_x\)，只有初始位置满足 \(x+v_xt\ge l\) 的分子在时刻 \(t\)
前撞过右壁。它们在管长中的初始位置所占比例是 \(\min(v_xt/l,1)\)。

\subsection{四、解题入口}\label{ux56dbux89e3ux9898ux5165ux53e3-5}

定义热速率 \[
v_T=\sqrt{\frac{2kT}{m}},
\] 并令 \[
u=\frac{v_x}{v_T},\qquad
a=\frac{l}{tv_T}. \tag{1}
\] 对 \(v_y,v_z\) 积分后，一维分布为 \[
\phi(u)\,du=\frac{1}{\sqrt\pi}e^{-u^2}\,du.
\] 气体总分子数由理想气体状态方程得到 \[
N=\frac{p_0Sl}{kT}. \tag{2}
\]

\subsection{五、解析}\label{ux4e94ux89e3ux6790-5}

\needspace{3\baselineskip}
\textbf{（一）累计撞击右壁的分子数}

给定 \(u>0\)。若 \(0<u<a\)，则 \(v_xt<l\)，只有靠近右壁、长度为 \(v_xt\)
的初始位置区间能撞壁，比例为 \[
\frac{v_xt}{l}=\frac ua.
\] 若 \(u\ge a\)，则任意初始位置都来得及撞壁，比例为 \(1\)。因此 \[
\frac{N_c}{N}
=\int_0^a\frac ua\frac{e^{-u^2}}{\sqrt\pi}\,du
+\int_a^\infty\frac{e^{-u^2}}{\sqrt\pi}\,du. \tag{3}
\] 第一项使用 \[
\int_0^a ue^{-u^2}\,du=\frac12(1-e^{-a^2}),
\] 第二项使用 \[
\int_a^\infty e^{-u^2}\,du
=\frac{\sqrt\pi}{2}-f(a).
\] 代入式（3）： \[
\boxed{
N_c=N\left[
\frac12-\frac{f(a)}{\sqrt\pi}
+\frac{1-e^{-a^2}}{2\sqrt\pi\,a}
\right]
}. \tag{4}
\]

检查两个极限：\(t\to0\) 时 \(a\to\infty\)，式（4）趋于
\(0\)；\(t\to\infty\) 时 \(a\to0\)，式（4）趋于 \(N/2\)，对应初始
\(v_x>0\) 的一半分子最终都撞过右壁。

\begin{figure}[H]
\centering
\includegraphics[width=0.82\linewidth,height=\textheight,keepaspectratio,alt={图6-2　初始位置---纵向速度平面中的撞壁区域}]{../figures/solution-06.pdf}
\caption*{图6-2　初始位置---纵向速度平面中的撞壁区域}
\end{figure}

\needspace{3\baselineskip}
\textbf{（二）传给右壁的总动量}

速度为 \(v_x>0\) 的分子完全弹性反射一次，传给右壁的正向动量为
\(2mv_x\)。因此 \[
\begin{aligned}
\Delta P
=\frac{2Nmv_T}{\sqrt\pi}\left[
\int_0^a u\frac ua e^{-u^2}\,du
+\int_a^\infty ue^{-u^2}\,du
\right].
\end{aligned} \tag{5}
\] 使用题给积分 \[
\int_0^a u^2e^{-u^2}\,du
=\frac12[f(a)-ae^{-a^2}],
\] 以及 \[
\int_a^\infty ue^{-u^2}\,du=\frac12e^{-a^2},
\] 式（5）中含 \(e^{-a^2}\) 的两项恰好抵消： \[
\begin{aligned}
\Delta P
&=\frac{2Nmv_T}{\sqrt\pi}
\left[
\frac{f(a)-ae^{-a^2}}{2a}
+\frac12e^{-a^2}
\right]\\
&=\boxed{\frac{Nmv_T}{\sqrt\pi\,a}f(a)}.
\end{aligned} \tag{6}
\] 由 \(1/a=tv_T/l\)、\(mv_T^2=2kT\) 和式（2），还可写成 \[
\boxed{
\Delta P=\frac{2p_0St}{\sqrt\pi}f(a)
}. \tag{7}
\]

\needspace{3\baselineskip}
\textbf{（三）自由试管的最终速度}

试管很重，故气体逸出过程中管的位移和速度对分布的一阶影响可忽略。初始总动量为零。最终所有分子都从左端离开：初始
\(v_x<0\) 的直接离开，初始 \(v_x>0\) 的撞右壁后以 \(-v_x\)
离开。因此最终每个分子的纵向动量都是 \(-m|v_x|\)。

一维麦克斯韦分布给出 \[
\begin{aligned}
\langle|v_x|\rangle
&=2v_T\int_0^\infty
u\frac{e^{-u^2}}{\sqrt\pi}\,du\\
&=\frac{v_T}{\sqrt\pi}.
\end{aligned}
\] 气体最终总动量为 \[
P_{\rm gas}=-\frac{Nmv_T}{\sqrt\pi}.
\] 由总动量守恒 \[
MV+P_{\rm gas}=0,
\] 所以 \[
\boxed{
V=\frac{Nmv_T}{M\sqrt\pi}
=\frac{p_0Sl}{M}
\sqrt{\frac{2m}{\pi kT}}
}. \tag{8}
\] 方向向右。

\needspace{3\baselineskip}
\textbf{（四）光子气体}

黑体辐射出射度与能量密度的关系为 \[
J=\frac{cu_\gamma}{4}.
\] 题给 \(J=\sigma T^4\)，所以 \[
u_\gamma=\frac{4\sigma T^4}{c}.
\] 管内初始光子总能量 \[
U_\gamma=u_\gamma Sl
=\frac{4\sigma T^4Sl}{c}. \tag{9}
\] 各向同性光子分布中 \[
\langle|\cos\theta|\rangle
=\frac{\int |\cos\theta|\,d\Omega}{\int d\Omega}
=\frac12.
\] 所有光子最终从左端逸出，最终纵向动量为 \[
P_{\gamma x}
=-\frac{U_\gamma}{c}\langle|\cos\theta|\rangle
=-\frac{U_\gamma}{2c}.
\] 由动量守恒： \[
MV_\gamma=\frac{U_\gamma}{2c}.
\] 代入式（9）： \[
\boxed{
V_\gamma=\frac{2\sigma T^4Sl}{Mc^2}
}. \tag{10}
\] 方向同样向右。

\textbf{撞壁计数与逸出计数的区别。} 原题左壁撤去、右壁保留，取原左端为
\(x=0\)，右壁为 \(x=l\)。只有初始 \(v_x>0\)
的分子可能撞右壁，其飞行距离为 \(l-x\)，所以时刻 \(t\) 前撞壁要求 \[
v_x>(l-x)/t.
\]
图6-2阴影在正速度阈值上方。对固定速度，还要按靠近右壁的初始位置比例加权；由此得到正文的
\(N_c\)，其长时极限为 \(N/2\)。

总逸出数是另一个量。初始向左分子的逸出距离为
\(x\)，向右分子先在右壁反射，逸出总路程为
\(2l-x\)。积分两个速度方向后，以 \(a=l/(v_Tt)\) 表示，得到 \[
\frac{N_{\rm out}}N=1-\frac{2f(2a)}{\sqrt\pi}+\frac{1-e^{-4a^2}}{2\sqrt\pi a}.
\]
它在长时间趋于一，与撞壁计数的极限不同。此式作为补充对照，不替换原题第一问。

右壁冲量只由撞壁分子贡献。单次完全弹性反射向右壁传递 \(2mv_x\)，因此 \[
J(t)=\frac Nl\int_0^l dx\int_{(l-x)/t}^{\infty}2mv\,f_v(v)\,dv,
\] 其中 \(f_v(v)=e^{-v^2/v_T^2}/(\sqrt\pi v_T)\)。长时积分为
\(Nm v_T/\sqrt\pi\)，与最终气体向左总动量的绝对值相等。

光子部分正文取只有初始腔内黑体辐射、撤壁后不继续补充的储量模型。若端壁被持续维持在温度
\(T\)
并不断发射，则累计反冲依赖时间，不能仍给同一个有限最终速度。热源约定必须与原题``热平衡''的解释同时说明。

\textbf{短时与长时撞壁数的两项极限。}
初始左、右速度分量各占一半。取时间很短，已经撞过右壁的正速度分子主要来自右壁附近；其余正速度分子需要更久才能撞壁，其贡献被高速度尾部压低。定义
\(a=l/(v_Tt)\)，短时对应 \(a\gg1\)。正文计数式展开为 \[
N_c/N\simeq\frac1{2\sqrt\pi a}
=\frac{v_Tt}{2\sqrt\pi l}.
\] 它正好等于初始右壁碰撞通量乘时间。

时间很长时 \(a\ll1\)，不能只把误差函数趋零而漏掉指数项的同阶修正。用 \[
\int_0^a e^{-u^2}du=a-a^3/3+O(a^5),
\] \[
1-e^{-a^2}=a^2-a^4/2+O(a^6),
\] 可逐项检查给定分组计数式的极限。撞壁数趋于
\(N/2\)，这是原题第一问的正确极限。若另求总逸出数，初始向左组按路程
\(x\)，向右组按反射后路程 \(2l-x\) 积分，最终总数才趋于 \(N\)。

全部分子最终逸出时，右壁只曾反射初始右行那一半，冲量仍为
\(Nm v_T/\sqrt\pi\)。左端出口带走的左向总动量与管的右向动量相抵，封闭的``管加全部分子''系统总动量保持初值零。

管壁若完全热化分子、使速度再次随机化，就不再是每分子仅有一次端板反射的模型。此时可能有多次碰撞，最终反冲需由边界散射与能量交换共同决定。题目的理想侧壁约定应在这些极限检查之前保留清楚。

\textbf{慢分子尾部使两个量以不同速率饱和。} 令 \(a=l/(v_Tt)\)，对长时间
\(a\ll1\) 展开正文的撞壁计数，得到 \[
N_c/N=\frac12-\frac{a}{2\sqrt\pi}+
\frac{a^3}{12\sqrt\pi}+O(a^5).
\] 还未撞壁的正速度分子比例首先按 \(t^{-1}\)
下降。它们主要是纵向速度很小的分子，不是远处持续流入的新分子。

端壁冲量的比值为 \[
\frac{J(t)}{J(\infty)}=\frac{f(a)}a
=1-\frac{a^2}{3}+\frac{a^4}{10}+O(a^6).
\] 因此冲量的相对缺额按 \(t^{-2}\)
下降。计数和冲量有不同的趋近幂次，因为尚未撞壁的慢分子每个携带的动量也较小。

总逸出数则为 \[
N_{\rm out}/N=1-\frac{2a}{\sqrt\pi}+
\frac{4a^3}{3\sqrt\pi}+O(a^5).
\]
这一系数与撞壁计数不同：先向右运动再反射的分子需要更长总路程。图6-1中左端开放、右端保留，才能与这些路径分组一致。

例如取 \(t=10l/v_T\)，即 \(a=0.1\)，首几项给出 \[
N_c/N\simeq0.47184,
\qquad J(t)/J(\infty)\simeq0.99668,
\] \[
N_{\rm out}/N\simeq0.88791.
\]
冲量已很接近最终值，但仍有约一成分子留在管内。不能用``动量接近饱和''推断``所有分子已逸出''。

短时间则撞壁数和冲量都先与时间成正比，其系数分别是初始分子数通量和压力乘面积。两个极限各有不同的物理含义，把同一个误差函数式用于两问时，需保留它前面的速度权重与归一化因子。

\newpage

\section{第 7 题}\label{ux7b2c-7-ux9898}

物屏、薄透镜 \(L_1\)、薄透镜 \(L_2\) 和眼睛依次共轴排列。物屏到 \(L_1\)
的距离为 \(l_1\)，两透镜间距为 \(l_2\)。已知 \[
2f_1=l_2,\qquad 2l_2=l_1,\qquad
D_1=3D_2,\qquad f_1=1.5f_2,
\] 且 \(D_1,D_2\ll f_1,f_2\)。

（1）判断哪个透镜是限制入眼光能的孔径光阑；

（2）求眼睛能看到的物屏最大半径 \(r\)；

（3）若物屏上某点发出的、通过第一透镜中心的光线还能进入眼睛，称该点``明亮''，求明亮区域半径
\(R\)。

\begin{figure}[H]
\centering
\includegraphics[width=0.82\linewidth,height=\textheight,keepaspectratio,alt={图7-1　物屏、双透镜与观察眼的轴向布局}]{../figures/question-07.pdf}
\caption*{图7-1　物屏、双透镜与观察眼的轴向布局}
\end{figure}

\subsection{一、题型定位}\label{ux4e00ux9898ux578bux5b9aux4f4d-6}

轴上孔径比较先确定光阑。离轴的最大可见范围只要求存在一条可通过的光线，允许擦过两镜边缘；题设的明亮范围另按第一镜中心光线判断。

\subsection{二、知识点}\label{ux4e8cux77e5ux8bc6ux70b9-6}

\begin{enumerate}
\def\labelenumi{\arabic{enumi}.}
\tightlist
\item
  近轴光线用 \((y,u)\) 表示；
\item
  自由传播 \(y'=y+du\)；
\item
  薄透镜折射 \(u'=u-y/f\)；
\item
  孔径光阑比较；
\item
  主光线与视场边缘；
\item
  给定通过透镜中心光线的几何判据。
\end{enumerate}

\subsection{三、物理图像}\label{ux4e09ux7269ux7406ux56feux50cf-6}

孔径光阑不是直径最小的透镜自动担任，而是要把不同孔径折算为物点处允许的角范围。第二问要在两镜允许的高度范围内寻找至少一条可通过的光线；只追光阑中心光线会漏掉渐晕区。第三问直接跟踪穿过
\(L_1\) 中心的指定光线。

\subsection{四、解题入口}\label{ux56dbux89e3ux9898ux5165ux53e3-6}

令 \[
f_2=f.
\] 由题给关系逐项得到 \[
f_1=\frac32f,\qquad
l_2=2f_1=3f,\qquad
l_1=2l_2=6f,\qquad
D_1=3D_2. \tag{1}
\] 以下令光线倾角的小量为 \(u\simeq\tan\alpha\)。

\subsection{五、解析}\label{ux4e94ux89e3ux6790-6}

\needspace{3\baselineskip}
\textbf{（一）轴上物点的孔径限制}

从轴上物点发出初始倾角为 \(u\) 的光线。到达 \(L_1\) 时高度 \[
y_1=l_1u=6fu. \tag{2}
\] 经过 \(L_1\) 后倾角变为 \[
\begin{aligned}
u_1
&=u-\frac{y_1}{f_1}\\
&=u-\frac{6fu}{(3/2)f}\\
&=u-4u=-3u.
\end{aligned} \tag{3}
\] 传播到 \(L_2\) 时高度 \[
\begin{aligned}
y_2
&=y_1+l_2u_1\\
&=6fu+(3f)(-3u)\\
&=-3fu.
\end{aligned} \tag{4}
\]

通过 \(L_1\) 要求 \[
|y_1|\le\frac{D_1}{2}
\quad\Longrightarrow\quad
|u|\le\frac{D_1}{12f}
=\frac{D_2}{4f}. \tag{5}
\] 通过 \(L_2\) 要求 \[
|y_2|\le\frac{D_2}{2}
\quad\Longrightarrow\quad
|u|\le\frac{D_2}{6f}. \tag{6}
\] 因为 \(D_2/(6f)<D_2/(4f)\)，所以限制更严格的是 \[
\boxed{L_2\text{ 为孔径光阑}}. \tag{7}
\]

\needspace{3\baselineskip}
\textbf{（二）最大可见范围：存在一条可通过的光线}

物点高度为 \(y_0\)，初始倾角为 \(u\)。两镜处高度满足 \[
y_1=y_0+6fu,\qquad y_2=-y_0-3fu.
\] 消去 \(u\) 得 \[
y_0=-y_1-2y_2. \tag{8}
\] 有限孔径要求
\(|y_1|\le D_1/2\)、\(|y_2|\le D_2/2\)。因此有光线通过的必要且充分条件是
\[
|y_0|\le\frac{D_1}{2}+D_2.
\] 临界光线取两镜同侧边缘
\(y_1=-D_1/2\)、\(y_2=-D_2/2\)；这是经折射连接两孔径的实际光线，不必经过光阑中心。于是
\[
\boxed{r_{\max}=\frac{D_1}{2}+D_2=\frac{5D_2}{2}}. \tag{11}
\]
等号是几何擦边界，有限光通量的可见点需在它内侧。这里沿用原图眼睛紧邻第二镜且不额外限光的理想模型；若眼瞳大小与距离另有规定，应增加眼瞳约束。

\begin{figure}[H]
\centering
\includegraphics[width=0.39\linewidth,height=\textheight,keepaspectratio,alt={图7-2　最大可见范围与指定中心光线的边界}]{../figures/solution-07.pdf}
\caption*{图7-2　最大可见范围与指定中心光线的边界}
\end{figure}

\needspace{3\baselineskip}
\textbf{（三）明亮区域}

设物屏上一点高度为 \(R\)。题目指定考虑通过 \(L_1\) 中心的光线，因此 \[
y_1=R+l_1u=0,
\] 即 \[
u=-\frac R{l_1}. \tag{12}
\] 光线通过 \(L_1\) 中心，不发生方向改变。到达 \(L_2\) 时 \[
\begin{aligned}
y_2
&=l_2u\\
&=-\frac{l_2}{l_1}R\\
&=-\frac R2.
\end{aligned}
\] 要进入 \(L_2\)，需 \[
\left|\frac R2\right|\le\frac{D_2}{2}.
\] 因此 \[
\boxed{R_{\max}=D_2}. \tag{13}
\]

\textbf{在两镜高度平面中寻找可行光线。} 令第二镜焦距为 \(f\)，第一镜为
\(3f/2\)，两段传播距离为 \(6f\) 和 \(3f\)。从物点 \((y_0,u)\)
出发，第一镜高度为 \[
y_1=y_0+6fu.
\] 第一镜后的斜率为 \(u_1=u-y_1/(3f/2)\)，传播到第二镜得到 \[
y_2=y_1+3fu_1=-y_0-3fu.
\] 消去 \(u\) 可得 \[
y_0=-y_1-2y_2.
\] 允许的 \((y_1,y_2)\)
构成矩形：\(|y_1|\le3D_2/2\)、\(|y_2|\le D_2/2\)。线性函数 \(-y_1-2y_2\)
的极值出现在矩形角点，因此 \[
|y_0|\le\frac32D_2+D_2=\frac52D_2.
\] 图7-2上图正采用同侧两边缘
\(y_1=-3D_2/2,y_2=-D_2/2\)。擦边极限只有零宽度的光束，实际有限亮度可见点应略在内侧。

若要求光线穿过第二镜中心，则 \(y_2=0\)，所得范围为
\(|y_0|\le3D_2/2\)。这是一个额外约束，漏掉了可见的渐晕区。若要求穿过第一镜中心，则
\(y_1=0\)，所得明亮边界为 \(|y_0|\le D_2\)，对应图中下图。

出第二镜后还会再次改变斜率： \[
u_2=u_1-y_2/f.
\] 上图擦边线的三段斜率分别为
\(-2D_2/(3f)\)、\(D_2/(3f)\)、\(5D_2/(6f)\)。这一数值校验保证折线的拐折方向与会聚透镜一致。若眼睛并非紧邻第二镜，需再传播至眼瞳平面并增加眼瞳高度限制。

\textbf{部分光束通过时的线性渐晕区。} 以 \(D_2\) 为高度单位、\(f_2\)
为长度单位，令 \(Y=y_0/D_2\)。两镜限制给初始斜率的两个区间： \[
\frac{-1.5-Y}{6}\le uf_2/D_2\le\frac{1.5-Y}{6},
\] \[
\frac{-.5-Y}{3}\le uf_2/D_2\le\frac{.5-Y}{3}.
\] 真正可通过的是两个区间的交集。对 \(Y\ge0\)，交集宽度按三段变化： \[
\Delta U=\begin{cases}
1/3,&0\le Y\le1/2,\\
(2.5-Y)/6,&1/2<Y<2.5,\\
0,&Y\ge2.5.
\end{cases}
\] 这里 \(\Delta U\)
是一维子午平面的角宽，不是二维光通量；二维光束面积还需对孔径圆盘的交集积分。

轴上时，第二镜允许角范围 \(\pm1/6\)，第一镜允许
\(\pm1/4\)，所以第二镜是孔径光阑。偏离轴后，第一镜开始截掉某些方向，允许角宽连续下降，直到
\(Y=2.5\) 只剩擦边一线。

题目指定穿第一镜中心的光线对应 \(U=-Y/6\)。它在第二镜高度为
\(-Y/2\)，所以只到 \(Y=1\)。这条特征光线的消失不等于所有方向都被截断：在
\(1<Y<2.5\)
仍有部分交集。若实际眼瞳远离第二镜，新增眼瞳约束会进一步削弱或改变这些范围，不能直接复用本段交集。

\textbf{二维圆孔径中的光束面积。}
子午平面的范围能给最大可见半径，但亮度还要看二维孔径。以 \(D_2\)
为单位，设物点横向矢量为 \(\bm Y\)，第二镜上的光线位置为
\(\bm z\)。线性传播逐分量满足 \[
\bm y_1=-\bm Y-2\bm z.
\] 两圆孔径给 \[
|\bm z|\le1/2,
\qquad |\bm z+\bm Y/2|\le3/4.
\] 所以允许光束是半径 \(r_a=1/2,r_b=3/4\) 的两个圆盘交集，圆心距为
\(d_c=|\bm Y|/2\)。交集非空要求 \(d_c\le r_a+r_b\)，正好恢复最大物点半径
\(5D_2/2\)。

第二镜圆盘完全包含在另一个圆盘内时，不发生由第一镜造成的孔径截光。这要求
\[
d_c\le r_b-r_a=1/4\quad\Longleftrightarrow\quad |\bm Y|\le1/2.
\] 因此``完整孔径通光''的半径为
\(D_2/2\)。原题所谓明亮边界按指定穿第一镜中心的光线定义为
\(D_2\)，两种判断采用不同标准；正文保留该题的特征光线约定，不把它等同于无渐晕范围。

在部分相交区，二维通光面积可由两个圆弓面积求得： \[
\begin{aligned}
A={}&r_a^2\arccos\frac{d_c^2+r_a^2-r_b^2}{2d_cr_a}
\\
&+r_b^2\arccos\frac{d_c^2+r_b^2-r_a^2}{2d_cr_b}
\\
&-\frac12\sqrt{H}.
\end{aligned}
\] \[
H=(-d_c+r_a+r_b)(d_c+r_a-r_b)
(d_c-r_a+r_b)(d_c+r_a+r_b).
\] 物点在轴上时 \(A=\pi/4\)；物点接近最大边界时
\(A\to0\)。二维亮度不能直接按前面一维角宽成正比，因为圆弓面积的变化不同。

这一面积比例还默认物点向各允许方向发射近似均匀、忽略余弦与透过率变化，且眼瞳不再截光。若要准确照度，还需使用实际光束权重。图7-2是几何边界图，能够证明是否有光线通过，不能单靠那一条擦边线计算有限亮度。

\textbf{第二镜后的光束方向。} 从物面到第二镜前，矩阵为 \[
\begin{pmatrix}y_2\\u_1\end{pmatrix}
=\begin{pmatrix}-1&-3f\\-2/(3f)&-3\end{pmatrix}
\begin{pmatrix}y_0\\u\end{pmatrix}.
\] 再经过焦距 \(f\) 的第二镜， \[
u_2=u_1-y_2/f=y_0/(3f).
\] 它与初始光线斜率 \(u\)
无关，所以同一物点通过的所有近轴光线在第二镜后互相平行。第一镜先在其右侧
\(2f\) 处形成中间实像，该像距离第二镜恰为 \(f\)，与上述矩阵结果一致。

眼睛看到的是随物点位置变化的平行光束方向，而不是第二镜后某个有限平面的像点。孔径是否允许某一光线通过，与出射光束是否平行是两个不同问题。

若眼瞳位于第二镜后距离 \(z_e\)，瞳面高度为 \[
y_e=y_2+\frac{z_e}{3f}y_0.
\] 因此原两镜孔径约束之外还要加入瞳面限制。取 \(z_e=3f/2\)，可用
\(y_0=-y_1-2y_2\) 化简为 \[
y_e=-y_1/2.
\]
这个平面是第一镜经第二镜成的孔径像平面，第一镜直径在此缩小为原来的一半。观察者位置与瞳孔大小会决定能否接收全部出射光束，不能只用两镜几何范围来保证实际眼睛看见同样范围。

图7-2下图第二镜后的直线段斜率为
\(y_0/(3f)\)，同一物点的其他通光光线也具有这个斜率。若在不同孔径位置任选两条光线，它们在第二镜后保持平行，能够直接检查最终像位在无穷远。改变眼瞳位置会改变收到的光线束，而不会改变这一步近轴折射所得的方向。
若眼瞳中心偏离光轴，须比较光线在瞳面的实际高度与瞳孔中心，而非只看物点到轴的距离。原题未给瞳孔直径和位置，故正文的两镜范围以观察者能接收出射光束为条件，不能再据此算出眼睛的绝对照度。

\newpage

\section{第 8 题}\label{ux7b2c-8-ux9898}

在地面惯性系 \(S\) 中，\(t=0\) 时火车左端 \(L\) 位于 \(x=0\)，右端 \(R\)
位于 \(x=l\)。两端从静止开始，始终具有相同的地面系坐标加速度
\(a\)，均向右运动，且只讨论 \(t<c/a\)。

（1）在地面时刻 \(t\)，求左端瞬时静止系中测得的瞬时加速度 \(a'\)；

（2）以左端本征时 \(\tau=0\) 对应 \(t=0,x=0\)，令
\(\sin\phi=at/c\)，写出左端世界线 \(c\tau(\phi)\) 和 \(x_L(\phi)\)；

（3）以右端本征时 \(\tau'=0\) 对应 \(t=0,x=l\)，写出右端世界线
\(c\tau'(\phi)\) 和 \(x_R(\phi)\)；

（4）地面时刻 \(t_A\) 的左端事件记为 \(A\)。在 \(A\)
对应的瞬时静止系中，求该时刻火车长度，并说明这种运动是否满足刚性运动所要求的相对论性长度关系。

\begin{figure}[H]
\centering
\includegraphics[width=0.8\linewidth,height=\textheight,keepaspectratio,alt={图8-1　两端具有相同坐标加速度的火车世界线}]{../figures/question-08.pdf}
\caption*{图8-1　两端具有相同坐标加速度的火车世界线}
\end{figure}

\subsection{一、题型定位}\label{ux4e00ux9898ux578bux5b9aux4f4d-7}

两端在地面系的距离固定，但左端瞬时静止系的同时线是倾斜的。必须先找到同时线与右端世界线的交点，再算长度，并核查交点仍在给定运动时段内。

\subsection{二、知识点}\label{ux4e8cux77e5ux8bc6ux70b9-7}

\begin{enumerate}
\def\labelenumi{\arabic{enumi}.}
\tightlist
\item
  一维速度变换与加速度变换；
\item
  本征时 \(d\tau=dt/\gamma\)；
\item
  参数代换 \(\beta=\sin\phi\)；
\item
  洛伦兹同时条件 \(\Delta t=v\Delta x/c^2\)；
\item
  瞬时静止系中的长度；
\item
  Born 刚性与坐标刚性的区别。
\end{enumerate}

\subsection{三、物理图像}\label{ux4e09ux7269ux7406ux56feux50cf-7}

地面系中两条世界线是相同抛物线的水平平移： \[
x_L(t)=\frac12at^2,\qquad
x_R(t)=l+\frac12at^2. \tag{1}
\] 地面同时线水平，沿它量得长度恒为 \(l\)。然而左端事件 \(A\)
的瞬时静止系同时线在 \(x-ct\) 图上倾斜，它与右端世界线相交于较晚的事件
\(B\)。真正的瞬时静止系长度必须由 \(A,B\) 两事件计算。

\subsection{四、解题入口}\label{ux56dbux89e3ux9898ux5165ux53e3-7}

两端地面系速度相同： \[
v=at,\qquad
\beta=\frac vc=\frac{at}{c},\qquad
\gamma=\frac1{\sqrt{1-\beta^2}}. \tag{2}
\] 一维平行加速度变换给出 \[
a'=\gamma^3a.
\] 本征时则从 \[
d\tau=dt\sqrt{1-\beta^2}
\] 积分。

\subsection{五、解析}\label{ux4e94ux89e3ux6790-7}

\needspace{3\baselineskip}
\textbf{（一）瞬时静止系加速度}

一维运动中，瞬时静止系的平行加速度与地面系坐标加速度关系为 \[
a'=\gamma^3a.
\] 代入式（2）： \[
\boxed{
a'=\frac{a}{[1-(at/c)^2]^{3/2}}
}. \tag{3}
\] 虽然地面系坐标加速度恒定，本征加速度并不恒定，并在 \(t\to c/a\)
时迅速增大。

\needspace{3\baselineskip}
\textbf{（二）左端的本征时参数方程}

令 \[
\sin\phi=\frac{at}{c}.
\] 则 \[
t=\frac ca\sin\phi,\qquad
dt=\frac ca\cos\phi\,d\phi,
\] 且在题目范围内 \(\cos\phi>0\)。于是 \[
\begin{aligned}
d\tau
&=dt\sqrt{1-\sin^2\phi}\\
&=\frac ca\cos^2\phi\,d\phi.
\end{aligned}
\] 从 \(0\) 积分到 \(\phi\)： \[
\begin{aligned}
\tau
&=\frac ca\int_0^\phi\cos^2\psi\,d\psi\\
&=\frac ca\int_0^\phi\frac{1+\cos2\psi}{2}\,d\psi\\
&=\frac c{2a}\left(\phi+\frac12\sin2\phi\right)\\
&=\frac c{2a}(\phi+\sin\phi\cos\phi).
\end{aligned}
\] 所以 \[
\boxed{
c\tau(\phi)=\frac{c^2}{2a}
(\phi+\sin\phi\cos\phi)
}. \tag{4}
\] 左端坐标为 \[
\boxed{
x_L(\phi)=\frac12a
\left(\frac ca\sin\phi\right)^2
=\frac{c^2}{2a}\sin^2\phi
}. \tag{5}
\]

\needspace{3\baselineskip}
\textbf{（三）右端世界线}

右端速度函数与左端完全相同，因此本征时积分相同： \[
\boxed{
c\tau'(\phi)=\frac{c^2}{2a}
(\phi+\sin\phi\cos\phi)
}. \tag{6}
\] 仅空间坐标平移 \(l\)： \[
\boxed{
x_R(\phi)=l+\frac{c^2}{2a}\sin^2\phi
}. \tag{7}
\] 这里两条世界线上使用同一数值的
\(\phi\)，表示同一地面时刻，不表示某个瞬时静止系中的同时。

\needspace{3\baselineskip}
\textbf{（四）瞬时静止系的同时事件}

令 \[
s=\frac{at_A}{c},\qquad
u=\frac{at_B}{c},\qquad
\lambda=\frac{al}{c^2},\qquad
\delta=u-s.
\] 事件 \(A\) 处左端速度为
\(v_A=cs\)。在以该速度运动的瞬时惯性系中，\(A,B\)
同时要求洛伦兹时间差为零： \[
\Delta t'=\gamma_A\left(\Delta t-\frac{v_A\Delta x}{c^2}\right)=0.
\] 所以 \[
\Delta t=\frac{v_A\Delta x}{c^2}. \tag{8}
\]

地面系中的两事件坐标差为 \[
\begin{aligned}
\Delta t&=t_B-t_A=\frac ca(u-s)=\frac ca\delta,\\
\Delta x
&=x_R(t_B)-x_L(t_A)\\
&=l+\frac a2(t_B^2-t_A^2)\\
&=\frac{c^2}{a}\left[
\lambda+\frac12(u^2-s^2)
\right].
\end{aligned} \tag{9}
\] 把式（9）代入式（8）： \[
\frac ca\delta
=\frac{cs}{c^2}\frac{c^2}{a}
\left[\lambda+\frac12(u^2-s^2)\right].
\] 约去 \(c/a\)，并用 \(u=s+\delta\)： \[
\begin{aligned}
\delta
&=s\left[\lambda+\frac12((s+\delta)^2-s^2)\right]\\
&=s\left(\lambda+s\delta+\frac12\delta^2\right).
\end{aligned}
\] 移项得到 \[
\frac s2\delta^2-(1-s^2)\delta+s\lambda=0. \tag{10}
\] 选择在 \(l\to0\) 时趋于 \(\delta\to0\) 的小根： \[
\boxed{
\delta=
\frac{(1-s^2)-
\sqrt{(1-s^2)^2-2s^2\lambda}}{s}
}\qquad(s>0). \tag{11}
\] 在 \(s=0\) 时连续极限为 \(\delta=0\)。

\begin{figure}[H]
\centering
\includegraphics[width=0.8\linewidth,height=\textheight,keepaspectratio,alt={图8-2　左端瞬时静止系同时线与右端世界线的交点}]{../figures/solution-08.pdf}
\caption*{图8-2　左端瞬时静止系同时线与右端世界线的交点}
\end{figure}

由式（8），洛伦兹变换后的空间差 \[
\Delta x'
=\gamma_A(\Delta x-v_A\Delta t)
=\gamma_A\Delta x(1-\beta_A^2)
=\frac{\Delta x}{\gamma_A}. \tag{12}
\] 又由同时条件 \[
\Delta x=\frac{c^2\Delta t}{v_A}
=\frac{c^2}{a}\frac{\delta}{s}.
\] 所以事件 \(A\) 的瞬时静止系测得火车长度为 \[
\boxed{
L_A=\frac{c^2}{a}\frac{\delta}{s}\sqrt{1-s^2}
}, \tag{13}
\] 其中 \(\delta\) 由式（11）给出。

当火车很短，满足 \[
\lambda=\frac{al}{c^2}\ll\frac{(1-s^2)^2}{s^2},
\] 对式（10）忽略 \(\delta^2\)，有 \[
\delta\simeq\frac{s\lambda}{1-s^2}.
\] 代入式（13）： \[
\boxed{
L_A\simeq\frac{l}{\sqrt{1-s^2}}
=\gamma_A l
}. \tag{14}
\] 这说明地面系长度恒定为 \(l\) 时，瞬时共动系中的长度随速度增大。它不是
Born
刚性运动。通常的长度收缩公式要求物体在自身静止系保持固定本征长度，而这里恰恰是本征长度在改变；因此并不与长度收缩矛盾。题面末句若称``钟慢效应''，从所求物理量看应理解为对相对论长度关系的检验，而不是对单个钟走时的比较。

\textbf{同时事件必须在轨迹定义域内。} 用无量纲变量
\(s=at_A/c\)、\(\lambda=al/c^2\)、\(\delta=a(t_B-t_A)/c\)。两事件的地面位置差为
\[
\frac{a\Delta x}{c^2}=\lambda+s\delta+\frac12\delta^2.
\] 瞬时静止系同时条件是 \(\Delta t=v_A\Delta x/c^2\)，因此 \[
\delta=s\left(\lambda+s\delta+\frac12\delta^2\right).
\] 选取随 \(l\to0\) 连续趋零的根， \[
\delta=\frac{1-s^2-\sqrt{(1-s^2)^2-2s^2\lambda}}s.
\] 判别式非负只保证形式交点存在。原题世界线只定义到 \(at/c<1\)，还必须有
\[
s+\delta<1.
\] 把边界 \(\delta=1-s\) 代回同时条件，得到有限定义域的限制 \[
\lambda<\frac{(1-s)^2(s+2)}{2s},\qquad0<s<1.
\]
若不满足这一条件，不能把超出时段的抛物线延长线当作原题已经给出的运动。

同时事件在瞬时静止系中的空间距离为 \[
L'=\gamma_A(\Delta x-v_A\Delta t)=\Delta x/\gamma_A.
\] 在 \(\lambda\) 很小时，\(\delta\simeq s\lambda/(1-s^2)\)，于是
\(L'\simeq\gamma_A l\)。这是短车近似，有限长度需用完整根。图8-2的 \(AB\)
连线已按这个同时条件绘制；水平连线则属于地面同时，二者测量的长度不同。

\textbf{一组有限长度事件的代回。} 取
\(s=0.5\)、\(\lambda=0.1\)，则判别式为 \(0.5125\)，连续小根为 \[
\delta=\frac{0.75-\sqrt{0.5125}}{0.5}\simeq0.06822.
\] 右端同时事件发生在 \(at_B/c=s+\delta\simeq0.56822<1\)。地面位置差为
\[
X=\lambda+s\delta+\delta^2/2\simeq0.13644.
\] 两事件满足 \(\delta=sX\)，可以直接代回洛伦兹同时式。左端洛伦兹因子为
\(\gamma_A=2/\sqrt3\)，故 \[
aL'/c^2=X/\gamma_A\simeq0.11816.
\] 短车近似
\(\gamma_A\lambda\simeq0.11547\)，已经有约百分之二的差异，说明有限车长公式有实际数值意义。

若保持 \(s=0.5\)、改取 \(\lambda=0.8\)，判别式仍为正，但有限时段上界为
\(0.625\)，所以同时交点超出已给定的世界线。应报告``运动时段内无所需同时事件''，而不是给出一个靠外推得到的火车长度。

在 \(s\to0\) 时根式有 \(0/0\)
外观。可先由同时条件展开，\(\delta\simeq s\lambda\)，于是
\(X\to\lambda\)、\(L'\to l\)，恢复静止长度。这个极限也检验了根的选择，另一大根不随静止极限连续。

\textbf{同一坐标加速度不等于同一固有加速度。}
原题两端在地面系保持相同加速度 \(a\)、同一时刻速度均为
\(at\)。对任一端，固有加速度为 \[
\alpha_{\rm prop}=\gamma^3a
=\frac{a}{(1-a^2t^2/c^2)^{3/2}}.
\]
它随地面时间增加，不是常量。因此不能套用恒定固有加速度的双曲线世界线替换题给抛物线。

在后端瞬时静止系中，与它同时的前端事件在更晚的地面时刻。原数值例子给
\(s_A=0.5,s_B\simeq0.56822\)，所以两同时事件上的速度也不同： \[
v_A=0.5c,
\qquad v_B\simeq0.56822c.
\] 变换前端速度得到 \[
v_B'=\frac{v_B-v_A}{1-v_Av_B/c^2}\simeq0.0953c.
\]
后端在此系中速度为零，前端仍在远离后端。于是有限长度瞬时静止系不是整节车厢共同的静止系，正文的
\(L'\) 是指定后端参考系中的同时事件间距。

短车近似的下一项也可直接由同时条件求出。令 \(\lambda\ll1\)，展开小根得
\[
\delta=\frac{s}{1-s^2}\lambda+
\frac{s^3}{2(1-s^2)^3}\lambda^2+O(\lambda^3).
\] 因为 \(X=\delta/s\)，所以 \[
\frac{aL'}{c^2}=\gamma_A\lambda+
\frac{s^2}{2(1-s^2)^{5/2}}\lambda^2+O(\lambda^3).
\] 二阶项为正，符合有限例子比 \(\gamma_A l\) 稍大的结果。接近 \(s=1\)
时展开系数迅速增大，需要更小的车长；``车长小''必须相对加速与速度尺度判断。

图8-2的斜线连接的两个事件既同时于后端参考系，又分属不同地面时刻。只画一条水平距离
\(l\)
再作尺缩，会漏掉前端在这段时间中继续前进的位移。这也是有限事件公式不能由单一静止尺的长度变换直接得到的原因。

\textbf{同时测长度时两端钟的读数。}
两端从地面同一时刻静止开始，具有相同坐标速度史；在地面同一时刻，两端积累的固有时间相同。但长度测量选取的
\(A,B\) 并不是地面同时事件。

对任一端，令 \(s=at/c\)，从初态到该事件的固有时间为 \[
\tau(s)=\frac ca\int_0^s\sqrt{1-u^2}\,du
=\frac{c}{2a}[\arcsin s+s\sqrt{1-s^2}].
\] 在原有限例子 \(s_A=0.5,s_B\simeq0.56822\) 中，前端钟读数较大： \[
\tau_B-\tau_A=\frac ca\int_{s_A}^{s_B}\sqrt{1-u^2}\,du>0.
\]
该差不是两钟在同一地面时刻走速不同，而是后端瞬时静止系的同时线截取了前端更晚的地面事件。

小车长极限中 \(\delta\simeq s\lambda/(1-s^2)\)，因此 \[
\frac{a(\tau_B-\tau_A)}c
\simeq\sqrt{1-s^2}\,\delta
=\frac{s\lambda}{\sqrt{1-s^2}}.
\] 长度的首项则为 \(aL'/c^2\simeq\lambda/\sqrt{1-s^2}\)，所以 \[
\tau_B-\tau_A\simeq v_A L'/c^2.
\] 这项同时性的钟差与长度来自同一对事件。若把前端钟读数强行取为
\(\tau_A\)，就换了一对事件，不能继续使用原同时距离。

图8-2的世界线和斜同时线同时表达位置与时间选择。读图时先确定参考系和事件，再计算钟或尺的量。两端钟虽在地面同时出发，读数相等的事件仍由共同的地面时间选取；后端瞬时系的同时测量则选取另一组事件。给出的钟差不是钟的故障，也不是额外的时间延迟，而是两种同时截面对同一世界线的不同截取。
\end{document}
